\documentclass[a4paper,12pt]{article}
\usepackage{amsmath,amsfonts,amssymb}
\usepackage{booktabs}
\usepackage{graphicx}
\usepackage{xcolor}
\usepackage{tikz}
\usepackage{pgfplots}
\pgfplotsset{compat=1.18}
\usetikzlibrary{arrows,positioning,calc,fadings,shapes,decorations.markings,decorations.pathmorphing,angles,quotes,patterns,arrows.meta}
\usepackage{circuitikz}
\usepackage{array}
\usepackage{multirow}
\usepackage{enumitem}
\usepackage{framed}
\usepackage{physics}
\usepackage{wrapfig}
\usepackage{caption}
\usepackage{empheq}
\usepackage[most]{tcolorbox}
\usepackage{pifont}
\usepackage{wasysym}
\usepackage{cancel}
\usepackage{xspace}
\usepackage{geometry}
\geometry{margin=2.5cm}
\usepackage[unicode,colorlinks]{hyperref}

% Colours
\colorlet{myred}{red!65!black}
\colorlet{xcol}{blue!70!black}
\colorlet{vcol}{green!60!black}
\colorlet{myblue}{blue!70!black}
\colorlet{mypurple}{blue!60!red!80}

% tcolorbox environments
\newenvironment{ilight}{%
  \centering\vspace*{4pt}%
  \begin{tcolorbox}[colframe=gray,colback=gray!10]%
}{%
  \end{tcolorbox}\vspace*{-4pt}%
}

\newenvironment{soln}{%
  \vspace*{4pt}%
  \begin{tcolorbox}[colframe=gray,colback=green!8,title=Solution]%
}{%
  \end{tcolorbox}\vspace*{-4pt}%
}

\newtcbox{\mymath}[1][]{%
  nobeforeafter,math upper,tcbox raise base,
  enhanced,colframe=blue!30!black,colback=blue!30,boxrule=1pt,#1}

% Counters
\newcounter{example}[section]
\newcounter{question}[section]
\renewcommand{\theexample}{\thesection.\arabic{example}}
\renewcommand{\thequestion}{\thesection.\arabic{question}}

\newcommand{\example}[1]{%
  \refstepcounter{example}%
  \noindent{\textcolor{cyan}{\textsf{\textbf{Example \theexample }}}\hspace*{1pt}#1}\vspace{4pt}}

\newcommand{\question}[1]{%
  \refstepcounter{question}%
  \noindent{\textcolor{magenta}{\textsf{\textbf{Question \thequestion }}}\hspace*{1pt}#1}}

% Useful macros from original
\newcommand{\keypoint}[1]{\textbf{\textcolor{red}{#1}}}
\newcommand{\cmt}{}
\newcommand{\sol}{\hspace{2cm}\textcolor{green!60!black}{\ding{52}}}
\newcommand{\RA}{\Rightarrow}
\newcommand{\ra}{\rightarrow}
\newcommand{\mps}{\ \text{m\,s}^{-1}}
\newcommand{\mpss}{\ \text{m\,s}^{-2}}
\newcommand{\jpkgk}{\text{J\,kg}^{-1}\text{K}^{-1}}
\newcommand{\NpC}{\text{N\,C}^{-1}}
\newcommand{\Vpm}{\text{V\,m}^{-1}}
\newcommand{\Vpcm}{\text{V\,cm}^{-1}}
\newcommand{\ec}{4\pi\epsilon_0}
\newcommand{\avg}[1]{\langle#1\rangle}
\newcommand{\eqyskip}{\vspace{4pt}}
\newcommand{\rcyskip}{\vspace{4pt}}
\newcommand{\OC}{^{\circ}\text{C}}
\newcommand{\xskip}{\hspace{1em}}
\newcommand{\tmin}{\text{min}}
\newcommand{\tmax}{\text{max}}
\newcommand{\fnet}{F_{\text{net}}}
\newcommand{\tcbhighmath}[1]{\boxed{#1}}
\newcommand{\radps}{\text{rad\,s}^{-1}}
\newcommand{\rad}{\text{rad}}
\newcommand{\muF}{\mu\text{F}}
\newcommand{\syl}[1]{}
\newcommand{\piste}{}
\newcolumntype{C}[1]{>{\centering\arraybackslash}p{#1}}
\newcolumntype{D}[1]{>{\centering\arraybackslash}m{#1}}

\tikzset{>=stealth'}
\tikzstyle{wave}=[blue!70!black,thick]

\begin{document}
\begin{center}
{\Huge\bfseries Worked Examples}\\[6pt]
{\large EMS Physics Notes}\\[4pt]
{\today}
\end{center}
\tableofcontents
\newpage
\section{Physical Quantities}

\example{Give the SI base units of (a) speed, (b) acceleration, (c) force, (d) work done.}

\begin{soln}
 $\text{speed} = \frac{\text{distance}}{\text{time}} \RA [v] = \frac{[s]}{[t]} = \frac{\text{m}}{\text{s}} = \text{m s}^{-1}$
	\end{soln}

\bigskip

\example{Alternative units of measurement for length}

\bigskip

\example{A ball falls in vacuum, all its gravitational potential energy converts into kinetic energy. This is expressed by the equation: $mgh = \frac{1}{2}mv^2$. Show that this equation is homogeneous.}

\begin{soln}
 LHS: $[mgh]=[m][g][h] = \text{kg} \times \mpss \times \text{m} = \text{kg m}^2\text{s}^{-2}$
	
	RHS: $\left[\frac{1}{2}mv^2\right] = [m][v]^2 = \text{kg} \times (\mps)^2 = \text{kg m}^2\text{s}^{-2}$
	
	so we see the equation $mgh = \frac{1}{2}mv^2$ is homogeneous \end{soln}

\bigskip

\example{The speed of a wave travelling along an elastic string is determined by three things: the tension $T$ in the string, the length $L$ of the string, and the mass $m$ of the string. Let's assume $v=T^a L^b m^c$, where $a$, $b$, $c$ are some numerical constants. Find the values of $a$, $b$ and $c$.}

\begin{soln} RHS: $[T]^a [L]^b [m]^c = (\text{kg m s}^{-2})^a \text{m}^b \text{kg}^c = \text{kg}^{a+c} \text{m}^{a+b} \text{s}^{-2a}$

for the equation to be homogeneous, we must have:
\begin{equation*}
	\text{kg}^{a+c} \text{m}^{a+b} \text{s}^{-2a} = \text{m s}^{-1} \RA
	\left\{ \begin{array}{ll}
	\text{kg}: & a+c=0 \\
	\text{m}: & a+b=1 \\
	\text{s}: & -2a=-1 \\
	\end{array}\right.
	\RA \left\{ \begin{array}{l}
	a=\tfrac{1}{2} \\
	b=\tfrac{1}{2} \\
	c=-\tfrac{1}{2} \\
	\end{array}\right.
\end{equation*}

so wave speed is given by: $v = T^{\frac{1}{2}} L^{\frac{1}{2}} m^{-\frac{1}{2}}$, or $v=\sqrt{\frac{TL}{m}}$

this happens to be the correct formula for the wave speed on a string \end{soln}

\bigskip

\example{A river flows from south to north with a speed of $2.0 \mps$ and the speed of a boat with respect to the water flow is $5.0 \mps$. (a) Suppose the boat leaves the west bank heading due east, what is the resultant velocity of the boat? (b) If the boat is to reach the exact opposite bank across the river, what is the resultant velocity and in what direction should the boat be headed?}

\begin{soln} Vector diagrams for the resultant velocity of the boat are illustrated below for part (a) boat heading due south:

magnitude of resultant velocity: $v = \sqrt{v_b^2 + v_w^2} = \sqrt{5.0^2 +2.0^2} \approx 5.4 \mps$.

In this case, the boat reaches opposite bank in shortest time but will drift downstream.
\end{soln}

\bigskip

\example{Given a vector $\overrightarrow{p}$, the graphical representations of $2\overrightarrow{p}$, $\frac{1}{2}\overrightarrow{p}$, $-\overrightarrow{p}$, $-\frac{3}{2}\overrightarrow{p}$ are:}

\bigskip

\example{A force of 3.0 N towards east and a force of 2.0 N towards 30$^\circ$ north of east act on an object. Find the magnitude and the direction of the resultant force.}

\begin{soln} suppose an arrow of length 1 cm represents a force of 1 N

one can draw a \emph{scale diagram} with a ruler and a protractor as shown
\end{soln}

\bigskip

\example{A box of weight $W=20.0 \text{ N}$ is resting on an inclined slope at 30$^\circ$ to the horizontal. Find the components of weight parallel to the slope and normal to the slope.}

\bigskip

\section{Kinematics}

\example{An athlete is running around a circular track of radius 60 m. When he completes one lap, what is the distance moved out? What about his displacement?}

\begin{soln}
    
Distance moved is the perimeter of the circle: $s = 2\pi r = 2\pi\times 60 \approx 380 \text{ m}$

athlete returns to same starting point after one lap, so displacement is zero
\end{soln}

\bigskip

\example{A ship travels 30 km north, takes a right, and then travels 40 km east to reach its destination. Compare the distance and the displacement travelled.}

\begin{soln}
Summing of all lengths gives distance: $30+40 = 70 \text{ km}$

The displacement vector is shown on the diagram.

$\text{The magnitude of displacement} = \sqrt{30^2 + 40^2} = 50 \text{ km}$

at an angle of $\theta = \tan^{-1}\frac{40}{30} \approx 53^\circ$ east of north
\end{soln}

\bigskip

\example{A cyclist travels a distance of 3.0 km in 20 minutes. She rests for 15 minutes. She then covers a further distance of 5.1 km in a time of 40 minutes. Calculate the average speed of the cyclist in m s$^{-1}$: (a) during the first 20 minutes of the journey, (b) for the whole journey.}

\begin{soln} For the first 20 minutes: $ v = \frac{3.0\times10^3}{20 \times 60} = 2.5 \mps$

\eqyskip For whole journey: $v = \frac{(3.0+0+5.1)\times10^3}{(20+15+40)\times60} = 1.8 \mps $ 
\end{soln}

\bigskip

\example{A man walks along a straight road for a distance of 800 m in 5.0 minutes. He then turns around, and walks along the same road for a distance of 280 m in 3.0 minutes. What is the average speed and the average velocity of this man during the 8.0 minutes?}

\begin{soln} $\text{Total distance travelled} = 800 + 280 = 1080 \text{m}$, so average speed: $v = \frac{1080}{8.0\times60} = 2.25 \mps$

\eqyskip  $\text{change of displacement} = 800 + (-280) = 520 \text{m}$, so average velocity: $v = \frac{520}{8.0\times60} \approx 1.08 \mps$ \end{soln}

\bigskip

\example{A maglev train travels at an average speed of $60\mps$ from the city centre to the airport, and at $40\mps$ on its return journey over the same distance. What is the average speed of the round-trip? What about the average velocity?}

\begin{soln} Suppose the distance between airport and city centre is $S$.

Average speed: $v = \frac{2S}{t_1 + t_2} = \frac{2S}{\frac{S}{60} + \frac{S}{40}} = 48 \mps$

For a round-trip, train returns to same staring position so change in displacement is zero. Average velocity is therefore \emph{Zero} \end{soln}

\bigskip

\example{A ball hits a barrier at right angles with a speed of $15 \mps$. It makes contact with the barrier for 30 ms and then rebounds with a speed of $12 \mps$. What is the average acceleration during the time of contact?}

\begin{soln} 
note that direction of velocity changed during rebound, so $\Delta v = 15 - (-12) = 27 \mps$

average acceleration: $a = \frac{\Delta v}{\Delta t} = \frac{27}{30 \times 10^{-3}} = 900 \mpss$ \end{soln}

\bigskip

\example{Describing the motion from the displacement-time graph shown.}

\bigskip

\example{The diagram shows the displacement-time graph for a vehicle travelling along a straight road. Use the graph to find, (a) the average velocity during the first 4.0 s of the motion, (b) the velocity of the vehicle \emph{at} time $t=1.5$ s.}

\begin{soln} 
during first 4.0 s, average velocity is

{
	\centering
	
	$ v = \frac{\Delta s}{\Delta t} = \frac{19.2}{4.0} \approx 4.8 \mps$
	
}

to find velocity at $t=1.5$ s, a tangent is drawn

gradient of tangent gives instantaneous velocity:

{
	\centering
	
	$ v = \frac{11.6 - 0}{4.0-0.75} \approx 3.6 \mps $
	
}
\end{soln}

\bigskip

\example{The velocity of a toy car is shown. For the journey shown on the graph, use the graph to find (a) the total distance travelled, and (b) the total displacement travelled.}

\begin{soln} Distance is estimated using area under $v$-$t$ graph

0$\sim$3.0 s: $s_1 = \frac{1}{2}\times{1.0+3.0}\times4.0 = 8.0 \text{ m}$

3.0$\sim$4.0 s: $s_2 = \frac{1}{2}\times{1.0}\times2.0 = 1.0 \text{ m}$

$\text{total distance} = 8.0 + 1.0 = 9.0 \text{ m}$

$\text{total displacement} = (+8.0) + (-1.0) = 7.0 \text{ m}$ 
\end{soln}

\bigskip

\example{Given the displacement-time graph as shown, check yourself that this $s$-$t$ graph leads to the velocity-time graph and the acceleration-time graph shown.}

\bigskip

\example{Given the velocity-time graph as shown, check yourself that this $v$-$t$ graph leads to the displacement-time graph as shown.}

\bigskip

\example{A car starts from rest and accelerates uniformly at 5.0 m s$^{-2}$ for 6.0 s.\\ (a) How fast is the car travelling at $t=8.0$ s? \\(b) What is the distance travelled by the car in this time?}

\begin{soln}
\begin{equation*}
v=u+at \RA v = 0 + 5.0\times 6.0 = 30 \mps
\end{equation*}
\begin{equation*}
s=ut+\frac{1}{2}at^2 \RA s= 0+\frac{1}{2}\times5.0\times6.0^2 = 90 \text{ m} 
\end{equation*}
\end{soln}

\bigskip

\example{A car is travelling at 30\mps. A hazard appears in front of the car, and the driver takes immediate action to stop the car. When brakes are applied, deceleration of the car is 5.0\mpss. What is the braking distance?}

\begin{soln}
\begin{equation*}
2as=v^2-u^2 \RA s = \frac{v^2-u^2}{2a} = \frac{0^2 - 30^2}{2\times(-5.0)} = 90 \text{ m} 
\end{equation*}
\end{soln}

\bigskip

\example{At the instant the traffic light turns green, a motorcycle waiting at the stop line starts with a constant acceleration of $2.0 \mpss$. At the same instant, a truck at a constant speed of $16 \mps$ overtakes and passes the motorcycle. How far beyond the stop line will the motorcycle overtake the truck?}

\begin{soln} Lets suppose overtake occurs at time $t$ after motorcycle starts to accelerate, so the distance travelled by motorcycle: $$s_m = u_mt + \frac{1}{2}at^2  \RA s_m = 0 + \frac{1}{2}\times 2.0 \times t^2$$

In the same time the distance travelled by truck: $$s_t = v_t t \RA s_t = 16t$$

The overtake when: $$s_m = s_t \RA \frac{1}{2}\times 2.0 \times t^2 = 16 t \RA t=16 \text{ s}$$

Substitute $t$ into either $s_m$ or $s_t$, one finds distance travelled: $s = 256 \text{ m}$ 
\end{soln}

\bigskip

\example{An object is released from rest from a height of $h=24$ m and falls freely under gravity. Air resistance is negligible.\\ (a) How long does it take to hit the ground? \\(b) What is its speed when hitting the ground?}

\begin{soln}
\begin{equation*}
h = \frac{1}{2}gt^2 \RA t = \sqrt{\frac{2h}{g}} = \sqrt{\frac{2\times24}{9.81}} \approx 2.21 \text{ s}
\end{equation*}
\begin{equation*}
v = gt = 9.81 \times 2.21 \RA v \approx 21.7 \mps
\end{equation*}

to find final velocity, one can also use $v^2 = u^2 + 2as$:
\begin{equation*}
v^2 = \cancelto{0}{u^2} + 2gh = 2\times9.81\times24 \RA v \approx 21.7 \mps 
\end{equation*}
\end{soln}

\bigskip

\example{A photograph is taken for a small particle falling from rest. The photograph is taken at 0.400 s after the object is released. Since the particle is still moving when the photograph is being taken, the image is blurred. It is found the blurred part corresponds to a length of 20.8 cm moved out by the particle, what is time of exposure for the photograph?}

\begin{soln} from $t=0$ to right before photo is taken: 
{	\centering
	
	$s_1 = \frac{1}{2}gt_1^2 = \frac{1}{2}\times9.81\times0.400^2 \approx 0.785 \text{ m}$
	
}

from $t=0$ to right after photo has been taken: 

{
	\centering
	
	$s_2 = s_1 + \Delta s = \frac{1}{2}gt^2_2 \RA t_2 = \sqrt{\frac{2(s_1+\Delta s)}{g}} = \sqrt{\frac{2\times(0.785+0.208)}{9.81}} \approx 0.450 \text{ s}$
	
}

time of exposure: $\Delta t = t_2 - t_1 = 0.450 - 0.400 \approx 0.050 \text{ s}$ 
\end{soln}

\bigskip

\example{A ball is projected vertically upwards at $12 \mps$. Air resistance is negligible. (a) Find the time taken for the ball to reach the highest position. (b) Find the greatest height.}

\bigskip

\example{A stone is thrown vertically upwards with an initial velocity of $14.0\mps$	from the edge of a cliff that is 35 m from the sea below.  \\(a) Find the speed at which it hits the sea. \\(b) Find the time taken for the stone to hit the sea.}

\bigskip

\example{A ball is thrown from a point $O$ at $15 \mps$ at an angle of $40^\circ$ to the horizontal. The ball reaches its highest position at point $P$. Ignore the effects of air resistance. (a) How long does it take to reach $P$? (b) What is the magnitude of the displacement $OP$?}

\bigskip

\example{A small object is horizontally projected at 7.20 ms$^{-1}$ from a surface at a height of $h=1.2 \text{ m}$ above the ground. Assume there is no air resistance. (a) What is the time taken for the object to hit the ground? (b) What is the horizontal range? (c) Find the velocity at which the object hits the ground.}

\begin{soln}
    
Vertically, take downward as positive: $h = \cancelto{0}{u_y t} + \frac{1}{2}gt^2 \ra t = \sqrt{\frac{2h}{g}} = \sqrt{\frac{2\times1.2}{9.81}} \approx 0.495 \text{ s} $

Horizontal range: $x = u_x t = 7.20 \times 0.495 \approx 3.56 \text{ m}$

Final vertical velocity: $v_y = \cancelto{0}{u_y} + gt = 9.81 \times 0.495 \approx 4.85 \mps$

Magnitude of resultant velocity: $v = \sqrt{v_x^2 + v_y^2} = \sqrt{7.20^2 + 4.85^2} \approx 8.68 \mps$

Angle to which resultant velocity makes with horizontal: $\phi = \tan^{-1}\frac{v_y}{v_x} = \tan^{-1}\frac{4.85}{7.20} \approx 34^\circ$ 
\end{soln}

\bigskip

\section{Force \\& Motion}

\example{A box of 6.0 kg is being pushed along a horizontal surface with a force of 30 N. The resistive force acting is 21 N. What is the acceleration of the box?}

\begin{soln}
\begin{equation*}
	\fnet = F - f = ma \RA a = \frac{F-f}{m} = \frac{30-21}{6.0} = 1.5 \mpss 
\end{equation*}
\end{soln}

\bigskip

\example{A car of mass 800 kg is travelling at a speed of $20 \mps$. The driver then operates the brake pedal so a braking force of 2000 N gradually brings the car to stop. (a) What is the deceleration for the car? (b) What is the stopping distance?}

\begin{soln} Using Newton's second law and noticing braking force acts opposite to direction of motion:

{
	\centering
	
	$\fnet = ma \RA -2000 = 800 \times a \RA a = -2.5 \mpss$
	
	$2as = v^2 - u^2 \RA s = \frac{v^2-u^2}{2a} = \frac{0^2 - 20^2}{2\times(-2.5)} = 80\text{ m}$
	
}
\end{soln}

\bigskip

\example{A massive ball is suspended on a string. A second string is attached to the bottom of the ball. If one pulls the bottom string with a gradually increasing force, does the top string or the bottom string break first? What if the bottom string is jerked, which string breaks?}

\bigskip

\example{Leviation of a slinky - When the top end of the slinky is dropped, the information of the tension change must propagate to the bottom end before both sides begin to fall; the top of an extended slinky will drop while the bottom initially remains in its original position, compressing the spring.}

\bigskip

\example{Suggest the action and reaction force in the following cases: (a) A person stands on a bathroom scale. (b) A helicopter hovers in air. (c) The earth orbits around the sun.}

\begin{soln}(a) The person exerts downward force on scale, scale exerts an upward reaction on the person.

(b) The rotors of helicopter push air downwards, air exerts an upward force on helicopter.

(c) The Sun pulls the earth through gravitational attraction, the Earth also attracts the sun in return. \end{soln}

\bigskip

\example{An astronaut finds that he weighs 300 N on the surface of Mars, where the gravitational field strength is known to be $3.7 \text{ N kg}^{-1}$. Find his mass and hence estimate his weight if he returns to his home on the Earth.}

\begin{soln} Mass of astronaut: $m = \frac{W_M}{g_M} = \frac{300}{3.7} \approx 81.1 \text{ kg}$

Weight on earth: $W_E = mg_E = 81.1 \times 9.81 \approx 795 \text{ N}$ \end{soln}

\bigskip

\example{An object of 5.0 kg falls through the atmosphere from a very high altitude. After some time, it falls at a constant speed of $70 \mps$. Assume there is no significant change in gravitational field during the fall and the air resistance is proportional to speed: $f=bv$. (a) Find the value of the coefficient $k$. (b) Find the acceleration of the object when it is falling at $30\mps$.}

\begin{soln}equilibrium between weight and air drag when falling at terminal speed, so
\begin{equation*}
	mg = bv_t \RA b = \frac{mg}{v_t} = \frac{5.0\times9.81}{70} \approx 0.70 \text{ kg s}^{-1}
\end{equation*}

at any instant, equation of motion is: $\fnet = ma = mg - bv$

at $30\mps$, acceleration is: $a = \frac{mg-bv}{m} = \frac{5.0\times9.81-0.70\times30}{5.0} \approx 5.6 \mpss$ \end{soln}

\bigskip

\example{A box of mass $m=4.0 \text{ kg}$ is resting on a horizontal ground. What is the normal contact force acting?}

\begin{soln} Equilibrium between weight and normal contact, so
\begin{equation*}
	R-W=0 \RA R=mg = 4.0 \times9.81 \approx 39.2 \text{ N} 
\end{equation*} \end{soln}

\bigskip

\example{A man of $80$ kg stands in a lift. Find his apparent weight, i.e., the contact force, when the lift is (a) moving upwards at steady speed of $2.0 \mps$, (b) accelerating upwards at $2.0 \mpss$, (c) moving upwards but slowing down at a deceleration of $1.5 \mpss$.}

\begin{soln} The forces acting on man are weight and normal contact, for either case, equation of motion for the man reads:

{

\centering

$\fnet = ma = R - mg$

}

so normal contact force: $R = mg + ma$

when rising at steady speed, man is in equilibrium ($a=0$), so: $R=mg = 80\times9.81 \approx 785 \text{ N}$

when accelerating upwards ($a = +2.0 \mpss$): $R = 80 \times 9.81 + 80 \times 2.0 \approx 945 \text{ N}$

when decelerating upwards ($a = -1.5 \mpss$): $R = 80 \times 9.81 + 80 \times (-1.5) \approx 665 \text{ N}$ \end{soln}

\bigskip

\example{A sleigh of mass 15 kg lies at rest on an icy ground. The surface is frictionless. A force $P$ of 75 N is applied to the sleigh. Find the normal contact force and the acceleration of the sleigh if $P$ is acting (a) horizontally, (b) at an angle $\alpha$ to the horizontal where $\tan\alpha=\frac{3}{4}$.}

\bigskip

\example{A block of mass $m$ stays at rest on an inclined plane. The plane makes an angle $\theta$ with the horizontal. Find the normal contact force $R$ and the frictional force $f$ acting on the block.}

\begin{soln} The block is in equilibrium, so $\fnet=0$ in any direction.

Parallel to slope: $f = W_\parallelslant \RA f=mg\sin\theta$

Normal to slope: $R = W_\perp \RA R =mg\cos\theta$ \end{soln}

\bigskip

\example{A block of mass $m$ slides down a \emph{smooth} slope. The angle of the slope to the horizontal is $\theta$. Find the acceleration of the block.}

\begin{soln} only force acting along the slope is component of weight down the slope, so:

{
	\centering
	
	$ \fnet = ma = mg\sin\theta \RA a = g\sin\theta $
	
}

as $\theta\to0$, $a\to0$, this shows if plane becomes horizontal, the block simply stays put

as $\theta\to90^\circ$, slope becomes vertical, block would undergo free fall, so naturally $a\to g$ \end{soln}

\bigskip

\example{A block of mass 2.0 kg slides down a rough slope from rest. The slope is inclined at angle $\theta=20^\circ$ to the horizontal, and the block experiences a constant friction of 5.0 N. (a) What is the block's acceleration? (b) What is the distance travelled in 2.5 seconds?}

\begin{soln} Resolving along slope:

{
	\centering
	
	$\fnet = mg\sin\theta - f = ma$
	
$a = \frac{2.0 \times 9.81\times\sin20^\circ - 5.0}{2.0} \approx 0.855 \mpss $\\
	
}

 Distance travelled: $s = \cancelto{0}{ut} + \frac{1}{2}at^2 = \frac{1}{2}\times0.855\times2.5^2 \approx 2.67 \text{ m}$ \end{soln}

\bigskip

\example{A block of mass 3.0 kg is travelling up an inclined slope at an initial speed of $2.8 \mps$. The slope makes an angle of $18^\circ$ with the horizontal. A constant friction of 7.5 N acts on the block.\\ (a) What is the block's deceleration?\\ (b) How far does the block travel along the slope before its speed decreases to zero? \\(c) Suggest whether the block could stay on the slope.\\}

\begin{soln} Resolving along slope (take direction of initial velocity as positive direction): 

{
	\centering
	
	$\fnet = -mg\sin\theta - f = ma \RA a = \frac{-mg\sin\theta-f}{m} = \frac{-3.0\times9.81\times\sin18^\circ-7.5}{3.0} \approx - 5.53 \mpss$
	
	\eqyskip $v^2 - u^2 = 2as \RA s = \frac{v^2 - u^2}{2a} = \frac{0^2 - 2.8^2}{2\times(-5.5.3)} \approx 0.709 \text{ m}$ 
	
}

Note that component of weight down the slope is: $W_\parallelslant = mg\sin\theta = 3.0\times9.81\times18^\circ \approx 9.1 \text{ N}$

$W_\parallelslant > f$, so friction is not enough to prevent block from sliding back down the slope \end{soln}

\bigskip

\example{Two boxes $A$ and $B$ are placed on a smooth surface. They are accelerated together by a horizontal force $F$ as shown. Find the acceleration and the contact force between them.}

\bigskip

\example{A vehicle of mass 1500 kg is towing a trailer of mass 500 kg by a light inextensible tow-bar. The engine of the vehicle exerts a driving force of 9600 N, and the tractor and the trailer experience resistances of 3600 N and 1800 N respectively. Find the acceleration of the vehicle and the tension in the tow-bar.}

\bigskip

\example{Two blocks of mass $m_A$ and $m_B$ ($m_A>m_B$) are joined together by a light inextensible string. The string passes over a smooth pulley as shown. The two blocks are suddenly released from rest. Find the acceleration of each block and the tension in the string.}

\begin{soln}
    
Apply Newton's second law to each block:

{
	\centering
	
	$\left\{ \begin{array}{llll} \text{for $A$:} & m_Ag-T &=& m_Aa \\ \text{for $B$:} & T-m_Bg &=& m_Ba \end{array} \right.$
	
}

\vspace{0.4em} adding the two, one obtains equation of motion for whole system:

{
	\centering
	
	$ m_A g - m_B g = (m_A+m_B) a $
	
}

		
solving these equations, we find
\begin{equation*}
	a = \frac{m_A-m_B}{m_A+m_B}g \qquad T = \frac{2m_Am_Bg}{m_A+m_B} 
\end{equation*}
\end{soln}

\bigskip

\example{A mass $M=4.0 \text{ kg}$ is attached to a block of mass $m=2.0 \text{ kg}$ through	a light string which passes over a frictionless pulley as shown. When both masses are released, find the acceleration and the tension in the string.}

\bigskip

\section{Mechanical Equilibrium}

\example{In the diagram, if we resolve the displacement; $$d_\perp = d\sin\theta$$ the moment of this force is: $$\tau = Fd\sin\theta$$\\
If we resolve the force: $$F_\perp = F\sin\theta$$ so moment of this force is: $$\tau = F\sin\theta d$$}

\bigskip

\example{The diagram shows a uniform beam of weight $W=20 \text{ N}$ and length 80 cm pivoted at point $P$. $P$ is 30 cm from one end. Two equal but opposite forces of magnitude $F = 12 \text{ N}$ are acting at the two ends of the beam as shown. What is the resultant moment about point $P$?}

\begin{soln} Moment of weight: $\tau_w = W d_w = 20 \times \left(0.50-\frac{1}{2}\times0.80\right) = 2.0 \text{ N m} \quad$ (clockwise)

Torque of couple: $\tau_c = F L = 12 \times 0.80 = 9.6 \text{ N m} \quad$ (anti-clockwise)

Resultant moment: $\tau_\text{net} = \tau_c - \tau_W = 9.6 - 2.0 = 7.6 \text{ N m} \quad$ (anti-clockwise) \end{soln}

\bigskip

\example{A uniform rod of length 90 cm is pivoted 30 cm from one end. It is balanced with a 3.0 kg load. Find the mass of the rod.}

\begin{soln}
    
 Take moments about $P$:
	
	$ 3.0 \times 9.81 \times 0.30 = m \times 9.81 \times \left(0.60-\frac{1}{2}\times0.90\right)$

so we find mass of rod: $m = 6.0 \text{ kg}$ 
\end{soln}

\bigskip

\example{A student balances a metre rule of mass 120 g supported on a fulcrum at the 40 cm mark. She then places a 20 g mass on the 70 cm mark and a 50 g mass on the 25 cm mark as shown. To balance the rule, what mass should she place on the 15 cm mark?}

\bigskip

\example{A cylinder of weight 100 N and diameter 50 cm rests against point $P$ of a curb of height 10 cm. What is the minimum force required to cause the cylinder to roll to the left?}

\bigskip

\example{A uniform plank of weight 100 N and length $L=6.0 \text{ m}$ rests horizontally on two supports $A$ and $B$. A man of weight 800 N stands a distance of $x=1.5 \text{ m}$ from end $A$. Determine the forces acting at the two supports.}

\begin{soln} Take moments about $A$: $ W_m x + W_p \cdot\frac{L}{2} = F_B L$

{
	\centering

	
	$ 800 \times 1.5 + 100 \times 3.0 = F_B \times 6.0 \ra F_B = 250 \text{ N} $
	
}

Take moments about $B$: $ W_m (L-x) + W_p \frac{L}{2} = F_A L$

{
	\centering

	$ 800 \times 4.5 + 100 \times 3.0 = F_A \times 6.0 \ra F_B = 650 \text{ N} $
	
}

One can check that: $F_A + F_B = W_m + W_p$, there must be no resultant force in vertical direction \end{soln}

\bigskip

\example{A uniform ladder of weight 120 N rests on a rough ground against a smooth wall as shown. The dimensions are labelled on the diagram. (a) What is the contact force acting at $B$? (b) What is the contact force acting at $A$? (c) What is the frictional force at $B$?}

\begin{soln} Free-body diagram is drawn as shown

resolve vertically: $R_B = W \RA R_B = 120 \text{ N}$

take moments about $B$: $R_A y = W \frac{x}{2}$

{
	\centering
	
	$ R_A = \frac{120 \times 0.30}{0.80} = 45 \text{ N}$
	
}

resolve horizontally: $f_B = R_A \RA f_B = 45 \text{ N}$ 
\end{soln}

\bigskip

\example{The diagram shows a uniform rod $AB$ of weight 60 N being held horizontally to a vertical wall by means of a light string. The string is attached to the rod at $B$, where a basket of weight 40 N is suspended. The other end of the string is fixed on the wall at $C$. The angle between the string and the rod is $30^\circ$. (a) Find the tension in the string. (b) Find the force acting on the rod at point $A$.}

\begin{soln} Take moments about $A$: $TL\sin\theta = W_b L + W_r  \frac{1}{2}L$

{
	\centering
	
	$ T \sin30^\circ = 40 + 60\times\frac{1}{2} \RA T=140 \text{ N} $
	
}

resolve horizontally: $F_{A,x} = T\cos\theta \RA F_{A_x}=140\cos30^\circ \approx 121 \text{ N}$

resolve vertically: $F_{A,y}+T\sin\theta = W_r + W_b \RA F_{A,y} = 60+40 - 140\sin30^\circ = 30 \text{ N} $

force at $A$: $F_A = \sqrt{F_{A,x}^2 + F_{A,y}^2} \RA F_A = \sqrt{121^2 + 30^2} \approx 125 \text{ N}$ \end{soln}

\bigskip

\example{A painting of weight $W=20\text{ N}$ is supported by two strings as shown. Both strings form an angle $\theta=30^\circ$ to the horizontal. Find the tension in the strings.}

\begin{soln} By resolving forces, we have: 

{
	\centering
	
	$\Biggl\{\begin{array}{l}
	T_1 \cos\theta = T_2 \cos\theta\\
	T_1\sin\theta + T_2\sin\theta = W
	\end{array} $ 
	
}

we solve the equations to obtain:

{
	\centering
	
	$ T_1 = T_2  = \frac{W}{2\sin\theta} = \frac{20}{2\sin30^\circ}=20\text{ N} $
	
}


alternatively, we can construct the force triangle as shown

$T_1$, $T_2$ and $W$ form an equilateral triangle, so
\begin{equation*}
	T_1=T_2=W=20\text{ N} 
\end{equation*}
\end{soln}

\bigskip

\example{The same painting of weight $W=20\text{ N}$ is supported by two strings at different angles $\theta_1=30^\circ$ and $\theta_2=45^\circ$ as shown. Find the forces in the two strings.}

\begin{soln} By resolving forces, we have: 

{
	\centering
	
	$ \left\{\begin{array}{l}
		T_1 \cos\theta_1 = T_2 \cos\theta_2\\
		T_1\sin\theta_1 + T_2\sin\theta_2 = W
	\end{array} \right.
	\RA
	\left\{\begin{array}{l}
	\frac{\sqrt{3}}{2}T_1 = \frac{\sqrt{2}}{2}T_2\\
	\frac{1}{2}T_1 + \frac{\sqrt{2}}{2}T_2 = 20
	\end{array} \right.
	$
	
}


solving this, we find: $T_1 \approx 14.6\text{ N}\, , \, T_2 \approx 17.9\text{ N} $

alternatively, we construct the force triangle as shown

$T_1$ and $T_2$  are related to $W$ by the the law of sine:

{
	\centering
	
	$\frac{W}{\sin75^\circ} = \frac{T_1}{\sin45^\circ} = \frac{T_2}{\sin60^\circ} $
	
}

from this we get the same result: $T_1 \approx 14.6\text{ N}\, , \, T_2 \approx 17.9\text{ N} $ \end{soln}

\bigskip

\section{Momentum}

\example{A ball of 120 g strikes a wall at right angle with a speed of $10 \mps$. It rebounds with the same speed. If the time of impact is 25 ms, find the average force exerted on the ball.}

\begin{soln} change in momentum: $\Delta p = mv - mu = 0.12 \times 10 - 0.12 \times (-10) = 2.4 \text{ kg m s}^{-1}$

average force: $F=\frac{\Delta p}{\Delta t} = \frac{2.4}{2.5 \times10^{-3}} = 96 \text{ N}$ \end{soln}

\bigskip

\example{Water is pumped through a hose-pipe. A man is holding the hose-pipe horizontally and water emerges from the hose-pipe with a speed of $16\mps$ at a rate of 45 kg per minute. Find the force required from this man to hold steady the hose-pipe.}

\begin{soln}
\begin{equation*}
	F = \frac{\Delta p}{\Delta t} = \frac{\Delta (mv)}{\Delta t} = \frac{\Delta m (v-0)}{\Delta t} = \frac{45 \times 16}{60} \RA F = 12\text{ N} 
\end{equation*}
\end{soln}

\bigskip

\example{A strong wind of speed $30\mps$ blows against a wall of area 10 m$^2$ at right angles. The density of the air is 1.2 kg m$^{-3}$. Assume air speed reduces to zero when it hits the wall. What is the approximate force exerted by the air on the wall?}

\begin{soln}

{
	\centering
	
	$ F = \frac{\Delta p}{\Delta t} = \frac{\Delta (mv)}{\Delta t} = \frac{\Delta m (v-0)}{\Delta t} \xlongequal{m=\rho V} \frac{\rho \Delta V v}{\Delta t} \xlongequal{V=AL} \frac{\rho A \Delta L v}{\Delta t} \xlongequal{\Delta L = v \Delta t} \rho A v^2 $
	\eqyskip\begin{equation*}
	\RA F = 1.2 \times 10 \times 30^2 = 10800 \text{ N} 
	\end{equation*}
	
}
\end{soln}

\bigskip

\example{Given the variation with time of the momentum of a body as shown in the $p$-$t$ graph, check yourself that the variation of force acting and the variation of the object's acceleration should be plotted as shown in the $F$-$t$ graph and the $a$-$t$ graph.}

\bigskip

\example{An object of mass 70 g is initially at rest. A force that varies with time is exerted on the object. The graph shows the how the force varies during the time of impact. What is the final velocity of the object?}

\begin{soln} area under $F$-$t$ graph gives change in momentum

{
	\centering
	
	$ \Delta p = \frac{1}{2} \times 50 \times 35\times10^{-3} = 0.875 \text{ kg m s}^{-1} $
	
}

final velocity: $v = \frac{\Delta p}{m} = \frac{0.875}{70\times10^{-3}} = 12.5 \mps $ \end{soln}

\bigskip

\example{A uranium-238 nucleus disintegrates, emitting an $\alpha$-particle of mass 4u and producing a thorium-234 nucleus of mass 234u. The uranium nucleus is initially at rest. (a) What is the ratio of the velocities of the product particles $\frac{v_\alpha}{v_\text{Th}}$?}

\bigskip

\example{A 3.0 kg mass moving at $6.0 \mps$ has a head-on collision with a stationary 1.0 kg mass. The two masses stick together on impact. What is the final velocity of the two masses?}

\begin{soln}\begin{equation*}
	m_1 u_1 + \cancelto{0}{m_2 u_2} = (m_1+m_2)v \RA 3.0\times6.0 = (3.0+1.0)\times v 
 \end{equation*}
 \begin{equation*}
 \RA v=4.5 \mps 
\end{equation*}
\end{soln}

\bigskip

\example{A 2.0 kg mass moving at $4.0 \mps$ collides head on with a 5.0 kg mass moving at $1.0 \mps$. After the collision, speed of the 5.0 kg mass is unchanged but its direction is reversed. What is the velocity of the 2.0 kg mass after the collision?}

\begin{soln}\begin{equation*}
m_1 u_1 + m_2 u_2 = m_1 v_1 + m_2 v_2 \, 
\end{equation*}
\begin{equation*}
\ra \, 2.0\times4.0+5.0\times(-1.0) = 2.0 \times v_1 + 5.0\times1.0 \, \ra \,  v_1 = -1.0 \mps
\end{equation*}

minus sign means the 2.0 kg mass reverses direction after collision \end{soln}

\bigskip

\example{A sphere of mass $m$ moves on a smooth horizontal surface at speed $v$ and collides \emph{elastically} with an identical ball at rest. What are the final velocities of the two spheres?}

\bigskip

\example{A 4.2 kg mass $A$ and a 1.5 kg mass $B$ are travelling towards each other on a frictionless horizontal plane. Mass $A$ and $B$ move at $5.40\mps$ and $1.80\mps$ respectively before they strike, as shown below. Mass $B$ moves to the right at $4.50\mps$ after the collision, (a) find the velocity of $A$ after the impact, and (b) suggest whether the collision is elastic.}

\bigskip

\example{A ball of mass $m_1 = 1.0 \text{ kg}$ travelling with a speed of $u=6.0 \mps$ in the $x$-direction strikes a stationary ball of mass $m_2 = 2.0 \text{ kg}$. The direction of the balls' velocities $v_1$ and $v_2$ after the collision are shown in the diagram. Find $v_1$ and $v_2$.}

\begin{soln}Start with momentum conservation equations:

\eqskip$\left\{\begin{array}{l}
	m_1 u = m_1 v_1 \cos60^\circ + m_2 v_2 \cos30^\circ \\
	0 = m_1 v_1 \sin60^\circ - m_2 v_2 \sin30^\circ
	\end{array} \right. $

\eqskip $ \left\{\begin{array}{l}
	1.0\times6.0 = 1.0 \times v_1 \times \frac{1}{2} + 2.0 \times v_2 \times \frac{\sqrt{3}}{2}  \\
	0 = 1.0 \times v_1 \times \frac{\sqrt{3}}{2} - 2.0 \times v_2 \times \frac{1}{2}
	\end{array} \right. $

\eqyskip Simplify and solve the equations:

\eqskip $ \left\{\begin{array}{l}
\frac{1}{2} v_1 + \sqrt{3} v_2 = 6 \\
v_2 =  \frac{\sqrt{3}}{2} v_1
\end{array} \right. \RA 
\left\{\begin{array}{l}
v_1 = 3.0 \mps \\
v_2 \approx 2.6 \mps
\end{array} \right.$

\eqyskip One can also draw and use the vector triangle for this question, this happens to be a right-angled triangle, so things become much easier.
\begin{equation*}
	\left\{\begin{array}{l}
	m_1 v_1 = m_1 u \cos 60^\circ \\
	m_2 v_2 = m_1 u \cos 30^\circ
	\end{array} \right. \RA 
	\left\{\begin{array}{l}
	1.0\times v_1 = 1.0 \times 6.0 \times \frac{1}{2} \\
	2.0\times v_2 = 1.0 \times 6.0 \times \frac{\sqrt{3}}{2}
	\end{array} \right.	\RA 
\end{equation*}
\begin{equation*}
	\left\{\begin{array}{l}
	v_1 = 3.0 \mps \\
	v_2 \approx 2.6 \mps
	\end{array} \right.  
\end{equation*}
\end{soln}

\bigskip

\section{Work \\& Energy}

\example{A 20 N force is applied at 60$^\circ$ to the horizontal to move a 1.0 kg object at a constant speed of 2.0$\mps$ for 30 s. How much work is done by the force?}

\begin{soln} \begin{equation*}
W = Fs\cos\theta = Fvt\cos\theta = 20 \times 2.0 \times 30 \times \cos60^\circ \RA W = 600\text{ J} 
\end{equation*}
\end{soln}

\bigskip

\example{A piston in a gas pump has an area of 600 cm$^2$. During one stroke, the pump moves a distance of 30 cm against a constant pressure of 8000 Pa. How much work is done?}

\begin{soln}
    
 \begin{equation*}
W = p \Delta V = p A \Delta s = 8000 \times 600 \times 10^{-4} \times 30 \times 10^{-2} \RA W = 144 \text{ J} 
\end{equation*}
\end{soln}

\bigskip

\example{When a spring is compressed by 2.0 cm, the force applied increases uniformly from zero to 4.0 N. How much work is done by this force?}

\begin{soln} $F$-$x$ graph for the force is plotted as shown

work to compress spring equals area under $F$-$x$ graph:
	$ W = \frac{1}{2} \times 4.0 \times 2.0 \times 10^{-2} = 0.040 \text{ J} $

 
 \end{soln}

\bigskip

\example{Estimate the kinetic energy of a running man.}

\begin{soln} Suppose the man has a mass of 75 kg and is running at $5 \mps$ (any reasonable value will do)
\begin{equation*}
	E_k = \frac{1}{2}mv^2 = \frac{1}{2}\times75\times5^2 \approx 940 \text{ J} 
\end{equation*}
\end{soln}

\bigskip

\example{A racing car of 800 kg starts off from rest. If the driving force is 5000 N, and the car experiences a constant resistive force of 1500 N, what is its speed after it has travelled 50 m?}

\begin{soln} The gain in K.E. equals work by driving force plus \emph{negative} work against resistance
\begin{equation*}
W_\text{total} = \Delta E_k \RA Fs - fs = \frac{1}{2}mv^2 - 0
\end{equation*}
\begin{equation*}
(5000 - 1500)\times50 = \frac{1}{2}\times800\times v^2 \RA v\approx20.9\mps 
\end{equation*}
\end{soln}

\bigskip

\example{A concrete cube of side 0.50 m and density 2400 kg m$^{-3}$ is lifted 4.0 m by a crane. How much work is done?}

\begin{soln}
    
The work done by crane equals gain in G.P.E. of cube
\begin{equation*}
W = \Delta E_p 
\end{equation*}
\begin{equation*}
\RA W = mg\Delta h = \rho V g \Delta h = 2400 \times 0.50^3 \times 9.81 \times 4.0 \approx 1.18\times10^4 \text{ J} 
\end{equation*}
\end{soln}

\bigskip

\example{For an object falling from rest due to gravity, if air resistance is negligible, what is its speed when it has fallen through a distance of $h$?}

\begin{soln}
\begin{equation*}
	\text{G.P.E loss }  =  \text{ K.E gain} \RA mgh = \frac{1}{2}mv^2 \RA v=\sqrt{2gh} 
\end{equation*}
\end{soln}

\bigskip

\example{A marble is projected vertically upwards with an initial velocity $u$. The average resistive force acting is $f$. How do you determine the maximum height reached by the marble?}

\begin{soln} $\text{K.E. loss} \,\,\, = \,\,\, \text{G.P.E. gain} \,\,\, + \,\,\, \text{energy loss due to resistance}$
\begin{equation*}
	\frac{1}{2}mu^2 - 0 = mgH_\tmax + f H_\tmax \RA  H_\tmax = \frac{mu^2}{2(mg+f)}
\end{equation*}
\end{soln}

\bigskip

\example{A box of mass $m$ slides down along a slope that is inclined at an angle $\theta$ to the horizontal. There is a constant friction $f$ acting on the box. When the box has moved through a distance of $\Delta s$ down the slope from $A$ to $B$, write down an equation relating its velocities $v_A$ and $v_B$ by applying the law of conservation of energy.}

\begin{soln}
    
\begin{eqnarray*}
		\text{Total work done} &=& \text{change in total energy} \\
		-W_f &=& \Delta E_k + \Delta E_p \\
		-fs &=& \left( \frac{1}{2}mv_B^2 - \frac{1}{2}mv_A^2 \right) + \left( mgh_B - mgh_A \right) \\
		-fs &=& \left( \frac{1}{2}mv_B^2 - \frac{1}{2}mv_A^2 \right) - mg\Delta h \\
		v_B^2 &=& v_A^2 + 2g\Delta s \sin\theta - \frac{2fs}{m}
	\end{eqnarray*}
	
	or equivalently, we can write
	\begin{eqnarray*}
		\text{G.P.E loss} &=& \text{K.E gain} \,\,\, + \,\,\, \text{energy loss due to friction}\\
		mg\Delta h &=& \left( \frac{1}{2}mv_B^2 - \frac{1}{2}mv_A^2 \right) + fs \\
		v_B^2 &=& v_A^2 + 2g\Delta s \sin\theta - \frac{2fs}{m}
	\end{eqnarray*}
	
	note that the two alternative ways of thinking produce the same result \end{soln}

\bigskip

\example{A slope is inclined at an angle $\theta$ to the horizontal. A box of mass $m$ is pushed up the slope with a constant force $F$ parallel to the slope, and the box experiences a constant frictional force $f$. When the box has moved through a distance of $\Delta s$ along the slope from $A$ to $B$, find an equation relating its velocities $v_A$ and $v_B$.}

\begin{soln}
\begin{eqnarray*}
	\text{Total work} &=& \text{change in total energy} \\
	W_F-W_f &=& \Delta E_k + \Delta E_p \\
	Fs-fs &=& \left( \frac{1}{2}mv_B^2 - \frac{1}{2}mv_A^2 \right) + \left( mgh_B - mgh_A \right) \\
	Fs-fs &=& \left( \frac{1}{2}mv_B^2 - \frac{1}{2}mv_A^2 \right) + mg\Delta h \\
	v_B^2 &=& v_A^2 + \frac{2(F-f)s}{m} - 2g\Delta s \sin\theta 
\end{eqnarray*}

or equivalently, we can write
\begin{eqnarray*}
	\text{work by $F$} &=& \text{K.E gain} \,\,\, + \,\,\, \text{G.P.E gain} \,\,\, + \,\,\, \text{energy loss due to friction}\\
	Fs &=& \left( \frac{1}{2}mv_B^2 - \frac{1}{2}mv_A^2 \right) + mg\Delta h + fs \\
	v_B^2 &=& v_A^2 + \frac{2(F-f)s}{m} - 2g\Delta s \sin\theta 
\end{eqnarray*}

again the two approaches produce the same expression for final velocity \end{soln}

\bigskip

\example{A \emph{ballistic pendulum} is a device used to measure the speeds of fast-moving bullets. It consists of a large block of wood of mass $M$, suspended from two long light strings. A bullet of mass $m$ is fired into the block, and the block and bullet combination swings upward. If the centre of mass rises a vertical distance $h$, what is the initial speed $u$ of the bullet?}

\begin{soln} As th bullet enters block, the combined momentum is conserved:
\begin{equation*}
	mu = (M+m)v \RA v = \frac{mu}{M+m}
\end{equation*}

when the system swings upward, K.E. transforms into G.P.E. but total energy is conserved:
\begin{equation*}
\frac{1}{2}(M+m)v^2 = (M+g)gh \RA v = \sqrt{2gh}
\end{equation*}

putting the two equations together, we find: $ u = \left(1+\frac{M}{m}\right)\sqrt{2gh} $ \end{soln}

\bigskip

\example{There are 150 steps to the top of a tower, and the average height of each step is 25 cm. It takes a man of 72 kg two minutes to run up all the steps. What is his average power?}

\begin{soln}
    
\begin{equation*}
	P = \frac{\Delta E_p}{\Delta t} = \frac{mg \Delta h}{\Delta t} = \frac{72\times9.81\times (150\times0.25)}{120} \RA P \approx 221 \text{ W} 
\end{equation*}
\end{soln}

\bigskip

\example{A turbine is used to generate electrical power from the wind. Given that the blades of the turbine sweep an area of 500 $\text{m}^2$, the density of air is	1.3 $\text{kg m}^{–3}$, and the wind speed is $10 \mps$. Assume no energy loss, find the power available from the wind.}

\begin{soln}\begin{equation*}
P = \frac{\Delta E_k}{\Delta t} = \frac{\frac{1}{2}\Delta m v^2}{\Delta t} = \frac{\frac{1}{2} \rho \Delta V v^2}{\Delta t} = \frac{\frac{1}{2} \rho A \Delta x v^2}{\Delta t} \RA P = \frac{1}{2} \rho A v^3 
\end{equation*}
\begin{equation*}
P = \frac{1}{2} \times 1.3 \times 500 \times 10^3 \approx 3.25 \times10^5 \text{ W}  
\end{equation*}
\end{soln}

\bigskip

\example{A ship is cruising a a constant speed of $15 \mps$. The total resistive force acting is 9000 N. What is the output power of this ship?}

\begin{soln} Constant speed so equilibrium between driving force and resistive force
\begin{equation*}
	P = F v = f v = 9000 \times 15 \RA P = 1.35 \times 10^5 \text{ W} 
\end{equation*}
\end{soln}

\bigskip

\example{A car of mass $800 kg$ accelerates from rest on a horizontal road. Suppose the engine provides a constant power of 24000 W, and the resistive force can be given by $f=16v$, where $v$ is the speed of the car in$\mps$. (a) What is the acceleration of the car when it is travelling at $15\mps$? (b) What happens to the car if it maintains this driving power?}

\begin{soln}The equation of motion for the car is: $ F_\text{net} = F - f = ma \RA \frac{P}{v} - \alpha v = ma $

at $v=15\mps$: $\frac{24000}{15} - 16\times15 = 800a \RA a=1.7\mpss $

as car's velocity $v$ increases, driving force $F=\frac{P}{v}$ decreases, resistive force $f=\alpha v$ increases

so resultant force will decrease, the car will accelerate at a decreasing rate

eventually it reaches an equilibrium state where $F=f$, the car then travels at constant speed $v_t$
\begin{equation*}
F=f \RA \frac{P}{v_t} = \alpha v_t \RA \frac{24000}{v_t} = 16v_t \RA v_t \approx 38.7\mps 
\end{equation*}
\end{soln}

\bigskip

\example{A water pumping system uses 3.0 kW of electrical power to raise water from a well. The pump lifts 1500 kg of water per minute through a vertical height of 8.0 m. What is the efficiency of the sysem?}

\begin{soln}
\begin{equation*}
	\eta = \frac{\Delta E_p}{\Delta E_\text{in}} = \frac{\Delta m g h}{P_\text{in} \Delta t} = \frac{1500\times9.81\times8.0}{3000\times60} \RA \eta \approx 65.4 \% 
\end{equation*}
\end{soln}

\bigskip

\example{Water flows into a turbine from a reservoir at a vertical distance of 70 m above. The water flows through the turbine at a rate of 2500 kg per minute. What is the output power of the turbine if it is 85\% efficient?}

\begin{soln}
\begin{equation*}
P_\text{out} = \eta P_\text{in} = \eta \frac{\Delta E_p}{\Delta t} = \eta \frac{\Delta m g h}{\Delta t} = 85\% \times \frac{2500 \times 9.81 \times 70}{60} 
\end{equation*}

\begin{equation*}
\RA P_\text{out} \approx 2.43 \times 10^4 \text{ W} 
\end{equation*}
\end{soln}

\bigskip

\section{Materials}

\example{A spring has a natural length of 20.0 cm. When a mass of 250 g is suspended from the spring, the new length of the spring is 26.0 cm. Find the spring constant.}

\begin{soln}\begin{equation*}
	k = \frac{F}{x} = \frac{mg}{L-L_0} = \frac{0.250\times9.81}{(26.0-20.0)\times10^{-2}} \RA k \approx 40.9 \text{ N m}^{-1} 
\end{equation*}
\end{soln}

\bigskip

\example{A spring has a spring constant of $270 \text{ N m}^{-1}$. A mass of 1.2 kg is hung from the spring. When the mass is released from a position where the spring has an extension of 5.0 cm, what is the acceleration of the mass?}

\begin{soln}\begin{equation*}
	\fnet = F - mg = kx - mg = ma \RA 
\end{equation*}
\begin{equation*}
 a = \frac{kx-mg}{m} = \frac{270\times5.0\times10^{-2} - 1.2 \times 9.81}{1.2} \approx 1.44 \mpss 
\end{equation*}
\end{soln}

\bigskip

\example{A steel spring has a spring constant of $20 \text{ N cm}^{-1}$. How much work is needed to stretch it from an extension of 3.0 cm to an extension of 5.0 cm?}

\begin{soln} Work done needed equals the increase in elastic potential energy:
\begin{equation*}
	W = \Delta E_p = \frac{1}{2}kx_2^2 - \frac{1}{2}kx_1^2 = \frac{1}{2} \times 2000 \times (0.050^2 - 0.030^2) = 1.6 \text{ J} 
\end{equation*}
\end{soln}

\bigskip

\example{A trolley of 400 g can travel freely along a horizontal surface. It is pushed against a spring buffer. Suppose the spring is initially compressed by 5.0 cm under a 20 N force. When the trolley is released, it accelerates until it becomes detached. What is the trolley's final speed?}

\begin{soln} Elastic potential energy in spring $\RA$ kinetic energy of trolley
\begin{equation*}
\frac{1}{2}Fx = \frac{1}{2}mv^2 \RA \frac{1}{2}\times20\times0.050 = \frac{1}{2}\times0.40\times v^2 \RA v\approx 1.58\mps 
\end{equation*}
\end{soln}

\bigskip

\example{The same spring is now set between two trolleys $A$ and $B$ of mass 400 g and 600 g. Initially the spring is again compressed by 5.0 cm under a force of 20 N. After both trolleys are released, what are their final speeds?}

\bigskip

\example{(a) A spring with $k_1 = 20 \text{ N cm}^{-1}$ is connected in series with a second spring with $k_2 = 30 \text{ N cm}^{-1}$. When a force of 60 N is applied, what is the total extension of the combination? (b) The same two springs are now connected in parallel. When a force of 50 N is applied on the combination, what is the extension?}

\begin{soln} For series connection: $k = \left(\frac{1}{k_1} + \frac{1}{k_1}\right)^{-1} = \left( \frac{1}{20} + \frac{1}{30} \right)^{-1} = 12 \text{ N cm}^{-1} \RA x = \frac{F}{k} = \frac{60}{12} = 5.0 \text{ cm}$

\eqskip or sum up extension of each spring: $x = x_1 + x_2 = \frac{F}{k_1} + \frac{F}{k_2} = \frac{60}{20} + \frac{60}{30} = 5.0 \text{ cm}$

for parallel connection: $k = k_1 + k_2 = 20 + 30 = 50 \text{ N cm}^{-1} \RA x = \frac{F}{k} = \frac{50}{50} = 1.0 \text{ cm} $ \eoe
\end{soln}

\bigskip

\example{A set of identical springs are set up as shown. Each individual spring extends by 1.0 cm under a load of 1.0 N. Assume the limit of proportionality is not exceeded, what is the total extension for this combination when a load of 6.0 N is applied?}

\begin{soln} For a single spring: $k = 1.0 \text{ N cm}^{-1}$, the combined spring constant: $k_\text{total} = \left( \frac{1}{3k} + \frac{1}{2k} + \frac{1}{k} \right)^{-1}$

{
	\centering
	
	\eqyskip $ k_\text{total} = \left(\frac{1}{3.0} + \frac{1}{2.0} + \frac{1}{1.0} \right)^{-1} = \frac{6}{11} \approx 0.545 \text{ N cm}^{-1}$
	
}

Total extension: $x = \frac{F}{k_\text{total}} = \frac{6.0}{0.545} = 11 \text{ cm}$

Alternatively, we can find and add extensions of each layer. In particular, the top layer withstands a force of 6.0 N shared by three springs, so each spring has a force of 2.0 N, extension is 2.0 cm. Similarly, the other two layers extend by 3.0 cm and 6.0 cm respectively, hence, total extension: $x = 2.0 + 3.0 + 6.0 = 11 \text{ cm}$ \end{soln}

\bigskip

\example{A copper wire has a cross-sectional area of $1.5 \times 10^{-6} \text{ m}^2$. The breaking stress of the wire is $2.0 \times 10^8 \text{ Pa}$. Find the breaking force.}

\begin{soln}
    
\begin{equation*}
\sigma = \frac{F}{A} \RA F = \sigma A = 2.0 \times 10^8 \times 1.5 \times 10^{-6} = 3.0 \times 10^2 \text{ N} 
\end{equation*}
\end{soln}

\bigskip

\example{A 200 N tensile force is applied on a steel wire of 1.5 m, the wire extends by 5.0 mm. The diameter of the cross section is 0.60 mm. What is the Young modulus of the steel wire?}

\begin{soln}\begin{equation*}
	E = \frac{\sfrac{F}{A}}{\sfrac{x}{L}} = \frac{FL}{Ax} = \frac{200 \times 1.5}{\pi \times (0.30\times10^{-3})^2 \times 5.0 \times 10^{-3}} \RA E \approx 2.1 \times 10^{11} \text{ Pa} 
\end{equation*}
\end{soln}

\bigskip

\example{A copper wire of length 2.0 m is under a stress of $7.8 \times 10^7 \text{ Pa}$. Given that the Young modulus of copper is $1.2 \times 10^{11} \text{ Pa}$, what is (a) the strain of the wire, (b) the extension of the wire?}

\begin{soln} strain: $\epsilon = \frac{\sigma}{E} = \frac{7.8 \times 10^7}{1.2 \times 10^{11}} \RA \epsilon = 6.5 \times 10^{-4} = 0.065 \%$

extension: $x = \epsilon L = 6.5 \times 10^{-4} \times 2.0 \RA x = 1.3 \times 10^{-3} \text{ m}$ \end{soln}

\bigskip

\example{Several blocks of steel are used to support a bridge. Each block has a height of 30 cm and a cross section of 15 cm $\times$ 15 cm. The steel block is designed to compress 2.0 mm when the maximum load is applied. Given that the Young modulus of steel is $2.1 \times 10^{11} \text{ Pa}$, what is the maximum load that can be supported by one block?}

\begin{soln} \begin{equation*}
	E = \frac{FL}{Ax} \RA F = \frac{EAx}{L} = \frac{2.1\times10^{11} \times (0.15\times0.15) \times2.0\times10^{-3}}{0.30} \end{equation*}
 
 \begin{equation*}
 \RA F \approx 3.15 \times 10^7 \text{ N} 
\end{equation*}
\end{soln}

\bigskip

\example{Two metal wires $A$ and $B$ are of the same length and they extend by the same amount under the same load. Given that Young modulus of wire $A$ is twice of $B$, what is the ratio of their diameters?}

\begin{soln}\begin{equation*}
	E = \frac{FL}{Ax} \RA A = \frac{1}{4}\pi d^2 = \frac{FL}{Ex} \RA 
 \end{equation*}
 \begin{equation*}
 d^2 \propto \frac{1}{E} \RA \frac{d_A}{d_B} = \sqrt{\frac{E_B}{{E_A}}} = \frac{1}{\sqrt{2}} 
\end{equation*}
\end{soln}

\bigskip

\example{A full-size crane is ten times greater than a model crane in all linear dimensions. If they are made of the same materials, what is the ratio of the cable's extension?}

\begin{soln} same material means same density $\rho$ and same Young's modulus $E$, so:
\begin{equation*}
E = \frac{FL}{Ax} \RA x = \frac{FL}{AE} = \frac{mgL}{\pi r^2 E} = \frac{\rho V g L}{\pi r^2 E} \RA
\end{equation*}
\begin{equation*}
x \propto \frac{VL}{r^2} \propto \frac{l^3 l}{l^2} \propto l^2 \RA \frac{x_\text{full}}{x_\text{model}} = 10^2 = 100 
\end{equation*}
\end{soln}

\bigskip

\section{Waves}

\example{When a wave travels on a water surface, the maximum depth of water is 21 cm and the minimum depth is 18 cm. What is the amplitude of the wave?}

\begin{soln} The amplitude is half the end-to-end distance: $ A = \frac{1}{2}(21-18) = 1.5 \text{ cm} $ \end{soln}

\bigskip

\example{A wave travelling at $4.0 \mps$ has a wavelength of 50 cm, what is its period?}

\begin{soln} \begin{equation*}
	v = \frac{\lambda}{T} \RA T = \frac{\lambda}{v} = \frac{0.50}{4.0} = 0.125 \text{ s} 
\end{equation*}
\end{soln}

\bigskip

\example{What is the distance between a compression and a rarefaction for a sound wave of frequency 550 Hz? (Speed of sound in air is about $330 \mps$.)}

\begin{soln}\begin{equation*}
	d = \frac{1}{2}\lambda = \frac{v}{2f} = \frac{330}{2\times 550} \RA d = 0.30 \text{ m} 
\end{equation*}
\end{soln}

\bigskip

\example{A sound wave is detected by a microphone and the trace is displayed on an oscilloscope as shown. If the time-base is set at 0.5 ms div$^{-1}$ and the voltage gain is set at 2 V div$^{-1}$. What is the frequency and the amplitude of the signal?}

\begin{soln} period: $T = 3 \times 0.5 \text{ ms} = 1.5 \times10^{-3} \text{ s}$

\eqyskip frequency: $f = \frac{1}{T} = \frac{1}{1.5 \times10^{-3} \text{ s}} \approx 667 \text{ Hz}$

\eqyskip amplitude: $A = 2 \times 2 \text{ V} = 4.0 \text{ V}$ \end{soln}

\bigskip

\example{A beam of electromagnetic radiation is known to have a frequency of 25 THz in vacuum. What is its wavelength?}

\begin{soln}
    
\begin{equation*}
	\lambda = \frac{c}{f} = \frac{3.00\times10^8}{25\times10^12} \RA \lambda = 1.2 \times 10^{-5} \text{ m} \quad \text{(infra-red)} 
\end{equation*}
\end{soln}

\bigskip

\example{If the amplitude of an incoming wave is increased by 50\%, what is the increase in the wave intensity?}

\begin{soln} $I \propto A^2 \RA \frac{I'}{I} = \left( \frac{A'}{A} \right)^2 = \left( 1 + 50\% \right)^2 =2.25 \RA $ so intensity is increased by 125\% \end{soln}

\bigskip

\example{Two observers $A$ and $B$ are at a distance $r_A$ and $r_B$ from a point source where $r_A=2r_B$. (a) Find the ratio of their intensities $\frac{I_A}{I_B}$. (b) Find the ratio of their amplitudes $\frac{A_A}{B_B}$.}

\begin{soln}

{
	\centering
	
	$ I \propto \frac{1}{r^2} \RA \frac{I_A}{I_B} = \left(\frac{r_B}{r_A}\right)^2 = \left(\frac{1}{2}\right)^2 \RA \frac{I_A}{I_B} = \frac{1}{4} $
	
	\vspace*{0.4em} $ I \propto A^2 \RA \frac{A_A}{A_B} = \sqrt{\frac{I_A}{I_B}} = \sqrt{\frac{1}{4}} \RA \frac{A_A}{A_B} = \frac{1}{2} $
	
}

\end{soln}

\bigskip

\example{A beam of light polarised in the vertical direction has an amplitude $A$ and intensity $I$. It passes through a polarising filter whose axis of polarisation is at $45^\circ$ to the vertical. (a) What is the amplitude and the direction of polarisation of the emerging beam? (b) If the emerging beam then enters another filter whose axis of polarisation is at $75^\circ$ to the vertical, what is the amplitude and intensity of the emerging beam?}

\begin{soln} through first filter: 
$A_1 = A \cos\theta_1 = A \cos 45^\circ \RA A_1 = \frac{1}{\sqrt{2}} A $
light is polarised in a direction at $60^\circ$ to the vertical

through second filter:
$A_2 = A_1 \cos\theta_2 = \frac{1}{\sqrt{2}} A \times \cos (75^\circ - 45^\circ) \RA A_2 = \frac{\sqrt{6}}{4} A $

$I_2 = \left(\frac{\sqrt{6}}{4} \right)^2 I \RA I_2 = \frac{3}{8} I$

 or, $I_2 = I \cos^2\theta_1 \cos^2 \theta_2^2 = I \times \cos^2 45^\circ \times \cos^2 (75^\circ-45^\circ) \RA I_2 = \frac{3}{8} I$ 

\end{soln}

\bigskip

\example{A police car moving towards you at $16 \mps$ sirens at 500 Hz. Given that the speed of sound in air is $340\mps$, at what frequency do you hear the siren?}

\begin{soln}
\begin{equation*}
	f' = f\frac{v}{v - v_s} = 500 \times \frac{340}{340-16} \RA f' \approx 525 \text{ Hz} 
\end{equation*}
    
\end{soln}

\bigskip

\example{A star emits an $H_\alpha$ line of wavelength 656 nm. (a) What is the frequency of this $H_\alpha$ line? (b) An observer on earth detects a wavelength of 680 nm, what can we say about motion of the star? (c) Find the relative speed of the star with respect to the earth.}

\begin{soln}
The original frequency of $H_\alpha$: $f = \frac{c}{\lambda} = \frac{3.00\times10^8}{656 \times 10^{-9}} \approx 4.57 \times 10^{14} \text{ Hz}$

The observed wavelength is longer (redshift), this means the star is moving away, or receding.

To find speed of star, we can consider observed frequency: $f' = \frac{c}{\lambda'} = f\frac{c}{c + v_s}$
\begin{equation*}
	\frac{3.00\times10^8}{680\times10^{-9}} = 4.57 \times 10^{14} \times \frac{3.00\times10^8}{3.00\times10^8 + v_s} \RA v_s \approx 1.10 \times 10^7 \mps
\end{equation*}

alternatively, we can consider observed wavelength: $\lambda' = \lambda + v_s T$, or: $\Delta \lambda = \lambda' - \lambda = \frac{v_s}{f}$
\begin{equation*}
	(680-656)\times10^{-9} = \frac{v_s}{4.57 \times 10^{14}} \RA  v_s \approx 1.10 \times 10^7 \mps 
\end{equation*}
\end{soln}

\bigskip

\section{superposition}

\example{Displacement–time graphs for two waves $A$ and $B$ meeting at some point, together with the resultant wave formed at that point.}

\bigskip

\example{Two loudspeakers are wired to produce identical sound signals in unison. The sound wave produced has a wavelength of 80 cm. Describe the volume you hear when you are at a distance of (a) 10 m from both speakers, (b) 10 m from one speaker and 12 m from the other?}

\begin{soln} path difference for case (a): $\Delta L = 10 - 10  = 0$

so constructive superposition, resultant amplitude is large, a loud sound is heard

path difference for case (b): $\Delta L = 12 - 10 = 2.0 \text{ m} \RA \Delta L = \frac{5}{2}\lambda$

so destructive superposition, resultant amplitude is small, sound is quiet \end{soln}

\bigskip

\example{The diagrams below each shows the displacement of a particular point in two transverse waves. If both waves are travelling in the same direction with a wavelength of 30 cm, what is the phase difference and path difference between each wave?}

\bigskip

\example{Two coherent waves meet in space, one of which has an amplitude of 0.30 cm and the other of 0.20 cm. How does the intensity of maxima compare with the intensity of minima?}

\begin{soln} Amplitude of maxima: $A_\tmax = A_1 + A_2 = 0.30 + 0.20 = 0.50 \text{ cm}$

Amplitude of minima: $A_\tmin = A_1 - A_2 = 0.30 - 0.20 = 0.10 \text{ cm}$

Ratio of intensities: $\frac{I_\tmax}{I_\tmin} = \left( \frac{A_\tmax}{A_\tmin} \right)^2 = \left( \frac{0.50}{0.10} \right)^2 \RA \frac{I_\tmax}{I_\tmin} = 25$ \end{soln}

\bigskip

\example{A teacher demonstrates the double-slit experiment with a beam of red light produced from a helium-neon laser. The beam of wavelength of 632 nm is sent through a double-slit separated by 0.30 mm and passed onto a wall at about 2.0 m from the slits. (a) What is the fringe separation? (b) If the experiment is carried out using a green laser, what changes do you expect?}

\begin{soln} Fringe separation: $ x = \frac{\lambda D}{d} = \frac{632 \times 10^{-9} \times 2.0 }{0.30 \times 10^{-3}} \RA x \approx 4.2 \times 10^{-3} \text{ m}$

if replaced by green light, wavelength becomes shorter, so smaller fringe separation \end{soln}

\bigskip

\example{Coherent light passes through a double slit. Initially, the light intensity from each slit is the same. State the change to the appearance of the fringes if (a) the separation between the slits is increased, (b) The width of one of the slits is reduced.}

\begin{soln} (a) From $ x = \frac{\lambda D}{d} $, shorter $\lambda$ results in smaller fringe separation

\hspace*{1.25em} but brightness of the fringes remains unchanged

(b) Reducing width of slit causes amplitude of light from that slit to decrease

\hspace*{1.25em} at maxima, resultant amplitude decreases, so bright fringes become less bright

\hspace*{1.25em} at minima, amplitudes do not cancel completely, so dark fringes become a bit brighter

\hspace*{1.25em} but no change to separation between fringes \end{soln}

\bigskip

\example{Light of 632 nm wavelength produces first-order maxima at 16$^\circ$ when passing through a diffraction grating at right angles. How many lines per millimetre are there in the diffraction grating?}

\begin{soln} slit separation: $d = \frac{n\lambda}{\sin\theta} = \frac{1\times632\times10^{-9}}{\sin 16^\circ} \RA d \approx 2.29 \times 10^{-6} \text{ m}$

number of lines in 1 mm: $N = \frac{1 \text{ mm}}{2.29 \times 10^{-6} \text{ m}} = \frac{1 \times 10^{-3}}{2.29 \times 10^{-6}} \RA N \approx 436$ \end{soln}

\bigskip

\example{Green light of 510 nm wavelength is incident normally on a diffraction grating with slit separation of 2.0 $\mu$m. (a) What is the highest order seen? (b) how many maxima are produced? (c) If blue light is used for the same diffraction grating experiment, suggest the effect on the diffraction pattern.}

\begin{soln} to find highest order, we have: $n_\tmax < \frac{d\sin 90^\circ}{\lambda} = \frac{2.0\times10^{-6}}{510\times10^{-9}} \approx 3.9 \RA n_\tmax = 3$

1st, 2nd, 3rd-order on either side and 0th-order at centre, so $ 2\times 3 + 1 = 7$ maxima are formed

blue light has shorter wavelength, so more maxima will be produced with blue light, and separation between each maxima would be closer \end{soln}

\bigskip

\example{Light of wavelength 630 nm produces a second-order maxima at angle of $60^\circ$ when it is directed at a diffraction grating. If visible light of another wavelength can also give a maximum at the same angle, what is this wavelength?}

\begin{soln} using diffraction grating equation: $d \sin\theta = n\lambda$, we find $n\lambda$ is constant for fixed $\theta$

from information given in the question, this constant is $ 2\times630 = 1260 \text{ nm}$

if $n=1$, $\lambda = 1260 \text{ nm}$, which would be infra-red

if $n=3$, $\lambda = 420 \text{ nm}$, which would be visible (violet)

if $n\geq 4$, $\lambda \leq 315 \text{ nm}$, which would be ultraviolet

so the desired wavelength is 420 nm \end{soln}

\bigskip

\example{\emph{Melde's experiment} demonstrates a stationary wave formed on a vibrating string using an oscillator set at a particular frequency. Given that the length of the string is 80 cm. (a) What is the wavelength of this wave? (b) If one wants to construct a stationary of longer wavelength, suggest what changes can be made?}

\bigskip

\example{The graph shows a standing sound wave produced in an air column in a closed pipe of length 30 cm. The frequency of this sound wave is 825 Hz. (a) What is the wavelength of the sound wave? (b) Calculate the speed of sound. (c) What is the frequency of the fundamental mode for sound wave in this air column? (d) Describe the motion of an air molecule near the open end.}

\bigskip

\example{\label{recpracwave}A tuning fork is held above a tall cylinder which is initially filled with water. The water level, measured from the bottom of the cylinder, is lowered at a constant rate. A loud sound is first heard when the water level is at 60.0 cm, and the next loud sound is heard when the water level is at 26.8 cm. It is known that the speed of sound in air is $340 \mps$. What is the frequency of the sound wave produced by the tuning fork?}

\begin{soln} loud sound is heard when stationary wave is set up

node is formed at surface of water

node-to-node distance is $60.0 - 26.8 = 33.2 \text{ cm}$

wavelength of this wave: $\lambda = 2 \times 33.2 \text{ cm} = 0.664 \text{ m}$

frequency: $f = \frac{v}{\lambda} = \frac{340}{0.664} \RA f \approx 512 \text{ Hz}$ 
\end{soln}

\bigskip

\section{Cosmology}

\example{The surface temperature of Proxima Centuri, the nearest star to Earth, is $3000$K and its luminosity is $6.506 \times 10^{23} $W.

Calculate the radius of Proxima Centuri in kilometres and show your working clearly.}

\begin{soln}
List the known quantities: 

Surface temperature, $T = 3000 K$\\
Luminosity, $L = 6.506 \times 10^{23} $W\\
Stefan's constant, $\sigma = 5.67 \times 10^{-8} W m^{-2} K^{-4}$\\
Write down Stefan's Law \& sub in values\\
The radius of Proximal Centuri is 106 200 km (4 s.f.)
\end{soln}

\bigskip

\example{The spectrum of the star Rigel in the constellation of Orion peaks at a wavelength of 263 nm, while the spectrum of the star Betelgeuse peaks at a wavelength of 828 nm.
Determine which of these two stars, Betelgeuse or Rigel, is cooler.}

\begin{soln}
List the known quantities
Maximum emission wavelength of Rigel = 263 nm = 263 × 10-9 m
Maximum emission wavelength of Betelgeuse = 828 × 10-9 m

$$T = \frac{2.9e-3}{\lambda_{max}}$$
Calculate the surface temperature of each star:
$$\text{Rigel: }T=\frac{2.9e{-3}}{263e{-9}}=11026 = 11 000K$$

$$\text{Betelgeuse: }T=\frac{2.9e{-3}}{828e{-9}}=3502 = 3500K$$

Step 4: Write a concluding sentence

Betelgeuse has a surface temperature of 3500 K, therefore, it is much cooler than Rigel
\end{soln}

\bigskip

\section{Electrical Quantities \& Components}

\example{A metal sphere carries a net charge of $+2.4 \times 10^{-9} \text{ C}$. Suggest whether the sphere has gained additional electrons or lost electrons, and also calculate how many electrons has the sphere gained or lost?}

\begin{soln} The sphere is positively-charged, so it has lost some of its electrons.

The number of lost electrons: $ N = \frac{Q}{e} = \frac{2.4\times10^{-9}}{1.60 \times10^{-19}} \RA N = 1.5 \times 10^{10}  $ \end{soln}

\bigskip

\example{A current of 25 mA flows in a wire. How many electrons have passed a point in the wire in one hour?}

\begin{soln}
    
Charge flow: $Q = I t = 25\times10^{-3} \times 3600 \RA Q = 90 \text{ C}$

Number of electrons: $N = \frac{Q}{e} = \frac{90}{1.60 \times10^{-19}} \RA N \approx 5.6 \times 10^{20}$ \end{soln}

\bigskip

\example{The current in a lamp is increased uniformly from 2.0 A to 8.0 A over 3.0 s. What is the charge flow?}

\begin{soln} Average current during this time is: $\bar{I} = \frac{2.0+8.0}{2} = 5.0 \text{ A} $

charge flow is then: $ Q = \bar{I} t = 5.0 \times 3.0 = 15 \text{ C} $

we can also use area under $I$-$t$ graph
\begin{equation*}
	Q = \frac{1}{2} \times (2.0 + 8.0) \times 3.0 = 15 \text{ C} 
\end{equation*}
\end{soln}

\bigskip

\example{A length of silver wire of a diameter of $1.2 \text{ mm}$ carries a current of 2.0 A. The electron number density in silver is $5.9\times10^{28} \text{ m}^{-3}$. What is the mean drift velocity of the electrons?}

\begin{soln}\begin{equation*}
v = \frac{I}{nAe} = \frac{I}{n\cdot\frac{1}{4} \pi d^2 \cdot e} \RA \end{equation*}
\begin{equation*}\frac{2.0}{5.9\times10^{28} \times \frac{1}{4}\pi \times (1.2\times10^{-3})^2\times1.60\times10^{-19}} \approx 1.9\times10^{-4} \mps
\end{equation*}

\eqyskip This calculation shows that each free electron actually travel at very low speeds in a wire \end{soln}

\bigskip

\example{A metal conductor and a piece of semi-conductor of the same cross section are connected in series. When an electric current is driven through them, compare the mean drift speed of the charge carriers in the two materials.}

\begin{soln}
    
The same current and same cross section, so mean drift speed $ v = \frac{I}{nAe} \propto \frac{1}{n}$

The metal has larger $n$, so free electrons move at relatively lower speeds

The semi-conductor has smaller $n$, so charge carriers move at relatively higher speeds \end{soln}

\bigskip

\example{How much electrical energy is transformed into thermal energy when a charge of 10 C flows through a heater whose p.d. across is 25 V?}

\begin{soln}\begin{equation*}
	V = \frac{W}{Q} \RA W = V Q = 25 \times 10 = 250 \text{ J} 
\end{equation*}
\end{soln}

\bigskip

\example{A fully-charged 12 V car battery can supply a total electrical energy of 0.90 MJ. The starter motor of the car requires an average current of 150 A for a period of 2.0s. The battery is not able to recharge due to a fault. How many times can the starter motor be used?}

\begin{soln} total charge that can be supplied: $Q_\text{total} = \frac{W}{\mathcal{E}} = \frac{0.90 \times 10^6}{12} = 75000 \text{ C}$

charge needed for each start: $Q = I t = 150 \times 2.0 = 300 \text{ C}$

number of starts possible: $N = \frac{Q_\text{total}}{Q} = \frac{75000}{300} = 250$ \end{soln}

\bigskip

\example{Electrons beams are accelerated between a heated cathode and an anode with a potential difference of 40 V. What is the speed of the electrons when they reach the anode?}

\begin{soln} electrical potential energy is transformed into kinetic energy of electrons:
\begin{equation*}
qV = \frac{1}{2}mv^2 \RA 
\end{equation*}
\begin{equation*}
v=\sqrt{\frac{2qV}{m}} = \sqrt{\frac{2 \times 1.60 \times 10^{-19} \times 40}{9.11 \times 10^{-31}}} \approx 3.75 \times 10^6 \mps 
\end{equation*}
\end{soln}

\bigskip

\example{A current of 0.10 A is driven through a heater when connected to a supply voltage of 220 V. (a) Find the resistance of the heater. (b) If the heater is made from a wire of resistivity $1.8\times10^{-5} \text{ }\Omega\text{ m}$ and a diameter of 0.24 mm, find the length of the wire.}

\begin{soln}\begin{equation*}
	R = \frac{V}{I} = \frac{220}{0.10} \RA R = 2200 \text{ }\Omega
\end{equation*}
\vspace*{-1.2em}\begin{equation*}
R = \frac{\rho L}{A} \RA L = \frac{RA}{\rho} = \frac{2200 \times \pi \times (0.12\times10^{-3})^2}{1.8 \times 10^{-5}} \RA L \approx 5.5 \text{ m} 
\end{equation*}
\end{soln}

\bigskip

\example{A uniform wire of resistance $R$ has a length of $L$ and a cross-sectional area of $A$. If the wire is stretched to twice the length while its volume remains constant, what is the new resistance $R'$ in terms of $R$?}

\begin{soln}volume: $V=LA = \text{constant}$, so $A \propto \frac{1}{L}$, hence doubling $L$ means $A$ would become halved

\vspace*{0.3em} resistance: $R = \frac{\rho L}{A} \propto \frac{L}{A} \RA \frac{R'}{R}  = \frac{L'}{L} \times \frac{A}{A'} = 2 \times 2 \RA R' = 4R $ \end{soln}

\bigskip

\example{An electric toaster is labelled `220 V 900 W'. When it is operating normally, find (a) the current through it, (b) the resistance in the toaster.}

\begin{soln}current: $I = \frac{P}{V} = \frac{900}{220} \RA I \approx 4.09 \text{ A}$

\vspace*{0.3em} resistance: $R = \frac{V}{I} = \frac{220}{4.09}$, $\,\,$ or alternatively, $R = \frac{V^2}{P} = \frac{220^2}{900} \RA R \approx 53.8 \text{ }\Omega$ \end{soln}

\bigskip

\example{Two conducting cylinders $X$ and $Y$ are of the same diameter and the same material. The length of the cylinder $X$ is twice of that of the cylinder $Y$. If the two cylinders are connected in series to a power supply so that the current through them are equal, find the ratio of the electrical powers dissipated in the two cylinders $\frac{P_X}{P_Y}$.}

\begin{soln} for either cylinder, electrical power: $P = I^2 R = I^2 \cdot \frac{\rho L}{A} = I^2 \cdot \frac{4 \rho L}{\pi d^2}$

same current $I$, same diameter $d$, same resistivity $\rho$, so power: $P \propto L \RA	\frac{P_X}{P_Y} = \frac{L_X}{L_Y} = 2 $ \end{soln}

\bigskip

\section{current}

\example{The diagram shows part of an electric circuit. The values of $I_1$, $I_2$ and $I_3$ are labelled. Calculate the current $I_4$.}

\begin{soln} applying the first law: $I_2 + I_3 = I_1 + I_4$
	
	\eqskip $\RA I_4 = 2.0 + 3.0 - 1.0 = 4.0$ A \end{soln}

\bigskip

\example{For the circuit shown, find an expression, in terms of $\mathcal{E}$ and $R$, for the electric current flowing through the resistors.}

\begin{soln} apply KVL for the closed loop shown:


	\centering
	
	$ \mathcal{E}_1 - \mathcal{E}_2 = IR_1 + IR_2 $
	
	\eqyskip $ I = \frac{\mathcal{E}_1 - \mathcal{E}_2}{R_1 + R_2} $
	
	\end{soln}

\bigskip

\example{The diagram shows a circuit where e.m.f. of the batteries and the values of resistance are all known. Write down the equations for the currents $I_1$, $I_2$ and $I_3$ using Kirchhoff's laws.}

\bigskip

\example{In the circuit shown, the battery has an e.m.f. $\mathcal{E}=18$ V, and the resistors have $R_1=6.0$ $\Omega$, $R_2=4.0$ $\Omega$ and $R_3=12$ $\Omega$. (a) Find the total resistance of all external resistors. (b) Find the current through the battery. (c) Find the p.d. across $R_1$, $R_2$ and $R_3$. (d) Find the total power dissipated in $R_2$ and $R_3$.}

\begin{soln} total resistance: $R = R_1 + \left( \frac{1}{R_2} + \frac{1}{R_2} \right)^{-1} = 6.0 + \left( \frac{1}{4.0} + \frac{1}{12} \right)^{-1} = 6.0 + 3.0 = 9.0 \text{ } \Omega$

\eqyskip current through battery: $I = \frac{\mathcal{E}}{R} = \frac{18}{9.0} = 2.0 \text{ A}$

p.d across $R_1$: $ V_1 = I R_1 = 2.0 \times 6.0 = 12 \text{ V}$

p.d across $R_2$ and $R_3$: $V_2 = V_3 = \mathcal{E} - V_1 = 18 - 12 = 6.0 \text{ V}$

power dissipated in $R_2$: $P_2 = \frac{V_2^2}{R_2} = \frac{6.0^2}{4.0} = 9.0 \text{ W}$

power dissipated in $R_3$: $P_3 = \frac{V_3^2}{R_3} = \frac{6.0^2}{12} = 3.0 \text{ W}$ \end{soln}

\bigskip

\example{A car battery has an e.m.f. of 20 V and an internal resistance of 0.50 $\Omega$. When a starter motor of resistance of 3.50 $\Omega$ is connected to the battery, find (a) the current supplied to the motor, (b) the p.d. across the battery terminals, (c) the power at which the motor operates.}

\begin{soln} current in circuit: $I = \frac{\mathcal{E}}{R+r} = \frac{20}{3.50 + 0.50} = 5.0 \text{ A} $

terminal p.d.: $V_R = IR = 5.0 \times 3.50 = 17.5 \text{ V}$

power of motor: $P_R = IV_R = 5.0 \times 17.5 = 87.5 \text{ W} \, \text{  or }\, P_R = I^2 R = 5.0^2 \times 3.50 = 87.5 \text{ W}$ \end{soln}

\bigskip

\example{A battery of e.m.f. $12$ V is connected to a network of total resistance of 22 $\Omega$. The current through the battery is $0.50$ A. Find (a) the internal resistance of the battery, (b) the total power produced by battery, (c) the power dissipated in battery.}

\begin{soln} lost volts: $V_r = \mathcal{E} - IR = 12 - 0.50 \times 22 = 1.0 \text{ V}$

internal resistance: $r = \frac{V_r}{I} = \frac{1.0}{0.50} = 2.0 \text{ }\Omega$

power produced by battery: $P_\text{total} = I\mathcal{E} = 0.50 \times 12 = 6.0 \text{ W}$

power dissipated in battery: $P_r = I V_r = 0.50 \times 1.0 = 0.50 \text{ W} \, \text{  or }\, P_r = I^2 r = 0.50^2 \times 2.0 = 0.50 \text{ W}$ \end{soln}

\bigskip

\example{An electric heater is operating at a working p.d. of 200 V. A current of 5.0 A is driven by a voltage source of 230 V. The source has an internal resistance of 4.0 $\Omega$ and the resistance of the connecting wires is not negligible. Find (a) the lost volts in the source, (b) the resistance of the wires, (c) the useful power from the source, (d) the efficiency of this circuit.}

\begin{soln} lost volts: $V_r = I r = 5.0 \times 4.0 = 20 \text{ V}$

p.d. across wires: $V_\text{wire} = \mathcal{E} - V_\text{heater} - V_r = 230 - 200 - 20 = 10 \text{ V}$

resistance of wires: $ R_\text{wire} = \frac{V_\text{wire}}{I} = \frac{10}{5.0} = 2.0 \text{ }\Omega$

useful power: $P_\text{useful} = P_\text{heater} = I V_\text{heater} = 5.0 \times 200  = 1000 \text{ W}$

efficiency: $\eta = \frac{P_\text{useful}}{P_\text{total}} = \frac{I V_\text{heater}}{I\mathcal{E}} = \frac{200}{230} \approx 87\%$ \end{soln}

\bigskip

\example{A cell of e.m.f. $\mathcal{E}$ and internal resistance $r$ is connected in series with a variable resistor $R$. $R$ is gradually increased from zero. (a) Suggest how the p.d. across the battery terminals change when $R$ is increased. (b) For larger values of $R$, power delivered to $R$ decreases. Suggest the advantage, despite the low power output, of using this cell in a circuit of larger resistance.}

\begin{soln} as $R$ increases, current in circuit decreases

less voltage is lost in cell, so terminal p.d. will increase

for same reason, less power is lost in cell, so higher efficiency for the circuit

more explicitly, terminal p.d. $V_R = I R = \frac{\mathcal{E}R}{R+r}$, and efficiency $ \eta = \frac{P_R}{P_\text{total}} = \frac{I^2 R}{I^2(R+r)} = \frac{R}{R+r}$

from these we can tell: if $R\up$, then $\frac{R}{R+r} \up$, so $V_R \up$ and $ \eta \up$ \end{soln}

\bigskip

\example{The diagram shows a simple circuit. The battery has a negligible internal resistance and an e.m.f. of 8.0 V. The resistor has a fixed resistance of 20 $\Omega$. (a) Find the readings shown on the ammeter and the voltmeter if both meters are ideal. (b) Instead, the ammeter has a non-zero resistance of 1.0 $\Omega$ and the voltmeter has a resistance of 60 $\Omega$, what are the true readings displayed on the two meters?}

\begin{soln} if both meters are ideal, then $I = \frac{\mathcal{E}}{R} = \frac{8.0}{20} = 0.40 \text{ A}$, and $V = \mathcal{E} = 8.0 \text{ V}$

for non-ideal case, total resistance in circuit: $R_\text{total} = R_A + \left( \frac{1}{R} + \frac{1}{R_V} \right)^{-1} = 1.0 + \left( \frac{1}{20} + \frac{1}{60} \right)^{-1} = 16 \text{ }\Omega$

current through ammeter: $I = \frac{\mathcal{E}}{R_\text{total}} = \frac{8.0}{16} = 0.50 \text{ A}$

p.d. across voltmeter: $V = \mathcal{E} - V_A = \mathcal{E} - I R_A = 8.0 - 0.50 \times 1.0 = 7.5 \text{ V} $ \end{soln}

\bigskip

\example{The diagram shows an electric circuit incorporating a power supply of negligible internal resistance and a high-resistance voltmeter. Determine the range of voltage reading on the voltmeter as the variable resistor $R$ is adjusted over its full range from 0 $\Omega$ to 50 $\Omega$.}

\begin{soln} $V_\tmin = \frac{R_\tmin}{R_\tmin + R_0} \times V_\text{in} = \frac{0}{0+10}\times 15 \RA V_\tmin = 0 \text{ V}$

\eqyskip $V_\tmax = \frac{R_\tmax}{R_\tmax + R_0} \times V_\text{in} = \frac{50}{50+10}\times 15 \RA V_\tmax = 12.5 \text{ V}$

so voltage reading ranges from 0 to 12.5 V \end{soln}

\bigskip

\example{Four resistors are connected to a battery as shown. (a) If the variable resistor $R_4$ is adjusted to have a resistance of 35 $\Omega$, what is the reading on the voltmeter? (b) If the voltmeter reads zero, what is the resistance for $R_4$?}

\begin{soln}
    
 (a) p.d. across $R_1$: $V_1 = \frac{10}{20+10}\times24 = 8.0 \text{ V}$

\eqyskip\hspace*{1.2em} p.d. across $R_3$: $V_3 = \frac{25}{35+25}\times24 = 10.0 \text{ V}$

\eqyskip\hspace*{1.2em} voltmeter reading: $V = V_3 - V_1 = 10.0 -8.0 = 2.0 \text{ V}$

(b) $V=0$ means p.d. share between $R_1$ and $R_2$ is same as p.d. share between $R_3$ and $R_4$
\begin{equation*}
	\frac{R_1}{R_2} = \frac{R_3}{R_4} \RA \frac{10}{20} = \frac{25}{R_4} \RA R_4 = 50 \text{ }\Omega 
\end{equation*}
\end{soln}

\bigskip

\example{In the circuit below, the driver cell has an e.m.f. of 17 V and an internal resistance of $5.0$ $\Omega$. $AB$ is a resistance wire of total resistance 80 $\Omega$ and length 80 cm. The e.m.f. of a test cell with an unknown internal resistance is to be determined. The moving contact $P$ is adjusted to a position where $AP=50$ cm such that the ammeter shows no reading. (a) Find the current flowing in the resistance wire. (b) Find the p.d across $A$ and $P$. (c) State the e.m.f. of the test cell.}

\bigskip

\section{\texorpdfstring{Electric Fields}

\example{An electron is at a point where the electric field strength is 1500 N C$^{-1}$, what is the electric force acting on the electron?}

\begin{soln}
    
\begin{equation*}
	F = Eq = 1500 \times 1.60 \times 10^{-19} \RA F = 2.4 \times 10^{-16} \text{ N} 
\end{equation*}
\end{soln}

\bigskip

\example{The diagrams below show the pattern of various electric fields.}

\bigskip

\example{The following examples show the set-up of two uniform electric fields.}

\bigskip

\example{Two horizontal parallel plate conductors are separated by 1.5 cm in air. A potential difference of 36 V is applied across the plates. Compare the electric force and gravitational force acting on a proton between the plates.}

\begin{soln}
    
 electric field strength: $E = \frac{V}{d} = \frac{36}{1.5\times10^{-2}} = 2400 \NpC$

electric force: $F_E = Eq = 2400 \times 1.60\times10^{-19} \approx 3.84\times10^{-16} \text{ N}$

weight of the proton: $W = mg = 1.67\times10^{-27} \times 9.81 \approx 1.64\times10^{-26} \text{ N}$

this shows electric force on proton are much stronger than gravitational forces

when dealing with motion of sub-atomic particles (e.g., protons, electrons, etc.) in an electric field, effects of gravity can therefore be ignored \end{soln}

\bigskip

\example{A uniform electric field is set up between the two parallel oppositely-charged metal plates separated by 8.0 cm. An oil droplet of mass 0.040 g and charge $-20$ nC is placed in the field. The droplet stays at rest. (a) Find the strength and the direction of the electric field. (b) What is the voltage required to produce this field? (c) If the separation between the plates is reduced, what would happen to the oil droplet?}

\begin{soln}(a) equilibrium so electric force equals weight: $F_E = mg = 0.050 \times 10^{-3} \times 9.81 \approx 3.92 \times 10^{-4} \text{ N}$

\hspace*{1.2em} field strength: $E = \frac{F_E}{q} = \frac{ 3.92 \times 10^{-4} }{ 20 \times 10^{-9}} \approx 1.96 \times 10^{4} \NpC$

\hspace*{1.2em} $F_E$ must act upwards for equilibrium, but droplet is negatively-charged

\hspace*{1.2em} so electric field acts in the downward direction

(b) voltage between the plates: $V = E d = 1.96 \times 10^{4} \times 8.0 \times 10^{-2} \approx 1.57 \times 10^3 \text{ V}$

(c) smaller separation means greater field strength, so greater electric force

\hspace*{1.2em} $F_E > mg$ so resultant force acts upwards, droplet will accelerate upwards \end{soln}

\bigskip

\example{A proton enters a uniform field of strength $1.0\times10^4\NpC$ at an initial velocity of $5.0\times10^5\mps$. The initial direction of motion is at right angles to the field. The dimensions of the field are shown in the diagram. (a) What is the time taken for the proton to pass through the field? (b) What is the deviation in the direction of the field during this time? (c) If an electron enters the field with the same initial velocity, describe how the deflection of the electron compares with that of the proton.}

\begin{soln} constant horizontal velocity, so time taken: $t = \frac{L}{u} = \frac{0.15}{5.0\times10^5} = 3.0\times10^{-7} \text{ s}$

\eqyskip constant acceleration in direction of field: $
a = \frac{F}{m} = \frac{Eq}{m} = \frac{1.0\times10^4\times1.6\times10^{-19}}{1.67\times10^{-27}} \approx 9.58\times10^{11}\mpss $

\eqyskip displacement moved in this direction: $
\Delta y=\frac{1}{2}at^2 = \frac{1}{2} \times 9.58\times10^{11} \times (3.0\times10^{-7})^2 \approx 0.043 \text{ m} $

electron would experience a much greater deflection in the opposite direction

\begin{itemize}

	\item[--] electron has opposite charge to proton, so deflect in the opposite direction
	
	\item[--] mass of electron is much smaller, so much greater acceleration
	
	time spent in the field is the same, so electron has much greater deflection 
\end{itemize}\end{soln}

\bigskip

\section{Sub-atomic Particles}

\example{The nuclide uranium-238 can be denoted by $^{238}_{\phantom{2}92}\text{U}$. (a) State its proton number and its neutron number. (b) Find the charge and mass of the uranium-238 nucleus.}

\begin{soln} proton number: $Z=92$

neutron number: $N=A-Z=238-92=146$

nuclear charge: $Q=+Ze=+92\times1.60\times10^{-19}\approx1.47\times10^{-17}\text{ C}$

mass of nucleus: $m=A\text{u} = 238 \times 1.66\times10^{-27} \approx 3.95 \times 10^{-25}\text{ kg}$ \end{soln}

\bigskip

\example{Radius of a copper-64 nucleus is $4.8\times10^{-15} \text{ m}$. Calculate the mean nuclear density.}

\begin{soln}
\begin{equation*}
	\rho = \frac{m}{V} = \frac{A\text{u}}{\frac{4}{3}\pi r^3} = \frac{64 \times 1.66\times10^{-27}}{\frac{4}{3}\pi \times (4.8x10^{-15})^3} \RA \rho \approx 2.3 \times 10^{17} \text{ kg m}^{-3}
\end{equation*}

this is much greater than density of a copper block (about $8.9 \times  10^3 \text{ kg m}^{-3}$)

we can therefore infer that the nucleus indeed takes most mass of the atom \end{soln}

\bigskip

\example{State and explain whether carbon-12 ($^{12}_{\phantom{1}6}\text{C}$) and carbon-14  ($^{14}_{\phantom{1}6}\text{C}$) are different isotopes of the same element.}

\begin{soln} both carbon-12 and carbon-14 atoms have 6 protons

but carbon-12 has 6 neutrons, carbon-14 has 8 neutrons

same proton number but different neutron number, so different isotopes of same element \end{soln}

\bigskip

\example{If a single beam of $\alpha$-particles, $\beta$-particles and $\gamma$-radiation is sent into a uniform electric field as shown, sketch and explain the paths of the beam.}

\bigskip

\example{The process in which a uranium-238 nucleus naturally decays into a thorium-234 nucleus through $\alpha$-emission can be written as:}

\bigskip

\example{Radon $^{222}_{\phantom{2}86}\text{Rn}$ decays in a sequence of processes to form bismuth $^{214}_{\phantom{2}83}\text{Bi}$ by emitting $\alpha$-particles and $\beta$-particles. For the decay chain of each radon nucleus, how many $\alpha$-particles and $\beta$-particles are emitted?}

\begin{soln} let $x$ be number of $\alpha$-particles emitted and $y$ be number of $\beta$-particles emitted
\begin{equation*}
	^{222}_{\phantom{2}86}\text{Rn} \longrightarrow ^{214}_{\phantom{2}83}\text{Bi} + x \cdot ^4_2 \alpha + y \cdot _{-1}^{\phantom{+}0}\beta + y \cdot _0^0 \bar{\nu}
\end{equation*}

Mass number and charge number are conserved:\\ 
$\Bigg\{$

\vspace*{-1.58\baselineskip}\hspace*{18pt} $\begin{array}{l}
4x + 214 = 222\\[-5pt]
2x - y + 83 = 86
\end{array} \RA x = 2$, $y=1$ \end{soln}

\bigskip

\example{Suggest the properties of the anti-particle of the electron.}

\begin{soln} anti-particle of electron has same electron mass: $m_e = 9.11\times10^{-31} \text{ kg}$

but it has a positive charge: $q = +e = + 1.60\times10^{-19} \text{ C}$

anti-particle of electron is usually called a \keypoint{positron} (positive electron) with the symbol $\text{e}^+$ \end{soln}

\bigskip

\example{By reference to the quark composition, explain the electric charge of protons and neutrons.}

\begin{soln} charge of proton (uud): $q_\text{p} = 2q_\text{u} + q_\text{d} = 2\times\left(+\frac{2}{3}e\right) + \left(-\frac{1}{3}e\right) = +e$

\eqyskip charge of neutron (udd): $q_\text{n} = q_\text{u} + 2q_\text{d} = \left(+\frac{2}{3}e\right) + 2\times\left(-\frac{1}{3}e\right) = 0$ \end{soln}

\bigskip

\example{A meson has an electric charge of $+e$ and is known to contain an up quark. Determine a possible flavour of the other quark.}

\begin{soln} charge of the other quark is: $(+e) - \left(+\frac{2}{3}e\right) = +\frac{1}{3}e$

this could be an anti-down ($\bar{\text{d}}$), an anti-strange ($\bar{\text{s}}$), or an anti-bottom quark ($\bar{\text{b}}$)\end{soln}

\bigskip

\example{Verify the conservation of (a) charge number, (b) mass number, and (c) lepton number for the $\beta^+$-decay process: $_1^1 \text{p} \longrightarrow _0^1 \text{n} + _{+1}^{\phantom{+}0}\beta + _0^0 \nu$}

\begin{soln} charge number conservation: $+1 = 0 + (+1) + 0$ \xskip \ding{51}

mass number conservation: $1 = 1 + 0 + 0$ \xskip \ding{51}

lepton number conservation: $0 = 0 + (-1) + (+1)$ \xskip \ding{51}

so this reaction does not violate any of these conservation law \end{soln}

\bigskip

\example{\emph{Electron capture} is a process where an electron in an atom's inner shell is drawn into the nucleus and combines with a proton. The result is to form a neutron and a neutrino. The process can be given by: $_1^1 \text{p} + _{-1}^{\phantom{+}0}\text{e} \longrightarrow _0^1 \text{n} +  + _0^0 \nu$. Show that the process satisfies the conservation of (a) charge number, (b) mass number, and (c) lepton number.}

\begin{soln} charge number conservation: $+1 + (-1) = 0 + 0$  \xskip \ding{51}

mass number conservation: $1 + 0 = 1 + 0$ \xskip \ding{51}

lepton number conservation: $0 + 1 = 0 + 1$  \xskip \ding{51}

so this reaction does not violate any of these conservation law \end{soln}

\bigskip

\section{Circular Motion}

\example{An object undergoes a uniform motion around a circular track of radius 2.5 m in 40 s, what is its angular speed and linear speed?}

\begin{soln}
    


\begin{equation*}
	\omega = \frac{2\pi}{T} = \frac{2\pi}{40} \approx 0.157 \radps \qquad v = \omega r = 0.157 \times 2.5 \approx 0.39 \mps 
\end{equation*}
\end{soln}

\bigskip

\example{Planet orbiting around a star}

\begin{soln}
\begin{center}
\begin{tikzpicture}[scale=0.8]
\shade [ball color = yellow] (0,0) circle (0.4);
\node at (0,-0.7){star};
\foreach \s in {0,30,60,...,330}
 \draw [orange, thick] (\s:0.5) -- ++(\s:0.1);
\draw [gray, thick, dashed] (0:2.5) arc [radius=2.5, start angle=0, end angle=360];
\shade [ball color = Cyan] (45:2.5) node[above right]{planet} circle (0.15);
\draw [purple, thick, ->] (45:2.5) -- ++(135:1.5) node[above]{$v$};
\draw [thick, ->, blue] (45:2.5) -- ++(225:1.35) node[above left]{gravity};
\end{tikzpicture}
\end{center}
gravity by the star provides centripetal force for the planet
\end{soln}

\bigskip

\example{A rock is able to orbit around the earth near the earth's surface. Let's ignore air resistance for this question, so the rock is acted by weight only. Given that radius of the earth $R=6400$ km. (a) What is the orbital speed of the rock? (b) What is the orbital period?}

\bigskip

\example{A turntable can rotate freely about a vertical axis through its centre. A small object is placed on the turntable at distance $d=40$ cm from the centre. The turntable is then set to rotate, and the angular speed of rotation is slowly increased. The coefficient of friction between the object and the turntable is $\mu = 0.30$. If the object does not slide off the turntable, find the maximum number of revolutions per minute.}

\begin{soln}if object stays on turntable, friction provides the centripetal force required: $f = m\omega^2 d$

increasing $\omega$ requires greater friction to provide centripetal force

but maximum limiting friction possible is: $f_\text{lim}  = \mu N = \mu mg$, therefore
\begin{equation*}
f \leq f_\text{lim} \RA m\omega^2d \leq \mu mg \RA \omega^2 \leq \frac{\mu g}{d} \end{equation*}
\begin{equation*}\RA \omega_\tmax = \sqrt{\frac{0.30\times9.81}{0.40}} \approx 2.71 \radps
\end{equation*}

period of revolution: $T_\tmin = \frac{2\pi}{\omega_\tmax} = \frac{2\pi}{2.71} \approx 2.32 \text{ s}$

number of revolutions in one minute: $n_\tmax = \frac{t}{T_\tmin} = \frac{60}{2.32} \approx 25.9 $ \end{soln}

\bigskip

\example{Particle $P$ of mass $m=0.40$ kg is attached to one end of a light inextensible string of length $r=0.80$ m. The particle is whirled at a constant angular speed $\omega$ in a vertical plane. (a) Given that the string never becomes slack, find the minimum value of $\omega$. (b) Given instead that the string will break if the tension is greater than 20 N, find the maximum value of $\omega$.}

\bigskip

\example{A pendulum bob of mass $120$ g moves at constant speed and traces out a circle or radius $r=10$ cm in a horizontal plane. The string makes an angle $\theta=25^\circ$ to the vertical. (a) What is the tension in the string? (b) At what speed is the bob moving?}

\bigskip

\example{A small ball of mass $m$ is attached to an inextensible string of length $l$.  The ball is held with the string taut and horizontal and is then released from rest.}

\bigskip

\section{Gravitational Fields}

\example{The Earth can be thought as a uniform sphere of radius $R = 6.4 \times 10^6$ m  and mass $M = 6.0 \times 10^{24}$ kg. Estimate the gravitational force on a man of 60 kg at sea level.}

\begin{soln} \begin{equation*}
F = \frac{GMm}{R^2} = \frac{6.67\times10^{-11}\times6.0\times10^{24}\times60}{(6.4\times10^6)^2} \approx 586 \text{ N} 
\end{equation*}
\end{soln}

\bigskip

\example{GPS (Global Positioning System) satellites move in a circular orbits at about 20000 km above the earth's surface. The Earth has a radius $R = 6.4 \times 10^6$ m  and mass $M = 6.0 \times 10^{24}$ kg. (a) Find the speed of GPS satellites. (b) Find its orbital period.}

\begin{soln}
\begin{equation*}
\frac{GMm}{r^2} = \frac{mv^2}{r} \RA v = \sqrt{\frac{6.67\times10^{-11}\times6.0\times10^{24}}{6.4\times10^6+2.0\times10^7}} \approx 3.9\times10^3 \mps
\end{equation*}
\begin{equation*}
v = \frac{2\pi r}{T} \RA T = \frac{2\pi r}{v} = \frac{2\pi\times(6.4\times10^6+2.0\times10^7)}{3.9\times10^3} \approx 4.3\times10^4 \text{ s}
\end{equation*}
\end{soln}

\bigskip

\example{A \keypoint{geostationary satellite} \index{geostationary orbit} moves in a circular orbit that appears motionless to ground observers. The satellite follows the Earth's rotation, so the satellite rotates from west to east above equator with an orbital period of 24 hours. Find the radius of this orbit.}

\begin{soln}
\begin{equation*}
\frac{GMm}{r^2} = m \omega^2 r \RA \frac{GMm}{r^2} = m\left(\frac{2\pi}{T}\right)^2 r \RA r^3 = \frac{GMT^2}{4\pi^2}
\end{equation*}
\begin{equation*}
r = \left( \frac{GMT^2}{4\pi^2} \right)^{1/3} = \left( \frac{6.67\times10^{-11} \times 6.0\times10^{24} \times (24\times3600)^2}{4\pi^2} \right)^{1/3} 
\end{equation*}
\begin{equation*}
    \approx 4.23\times10^7 \text{ m} 
\end{equation*}
\end{soln}

\bigskip

\example{Assuming the planets in the solar system all move around the sun in circular orbits, show that the square of  orbital period is proportional to the cube of orbital radius.
	\footnote{This is known as \emph{Kepler's 3rd law} for planetary motions. In the early 17th century, German astronomer Johannes Kepler discovered three scientific laws which describes how planets move around the sun. This $T^2 \propto r^3$ relation not only holds for circular orbits but are also correct for elliptical orbits.
		
		Isaac Newton proved that Kepler's laws are consequences of his own law of universal gravitation, and therefore explained why the planets move in this way.}}

\begin{soln}
    
\begin{equation*}
\frac{GMm}{r^2} = m \omega^2 r \RA \frac{GMm}{r^2} = m\left(\frac{2\pi}{T}\right)^2 r \RA T^2 = \frac{4\pi^2}{GM}\cdot r^3
\end{equation*}

$G$ is gravitational constant, $M$ is mass of the sun, so $\frac{4\pi^2}{GM}$ is a constant, so $T^2 \propto r^3$ \end{soln}

\bigskip

\example{A stone of mass 5.0 kg is hung from a newton-meter near the equator. The Earth can be considered to be a uniform sphere of radius $R = 6370$ km  and mass $M = 5.97 \times 10^{24}$ kg. (a) What is the gravitational force on the stone? (b) What is the reading on the meter?}

\begin{soln} gravitational force: $F_\text{grav} = \frac{GMm}{R^2} = \frac{6.67\times10^{-11}\times5.97\times10^{24}\times5.0}{(6.37\times10^6)^2} \approx 49.07 \text{ N}$

centripetal force required: $F_c = m\omega^2 R = m\left(\frac{2\pi}{T}\right)^2R = 5.0\times\frac{4\pi^2}{(24\times3600)^2}\times6.37\times10^6 \approx 0.17 \text{ N}$

apparent weight, or reading on meter: $N = F_\text{grav} - F_c = 49.07 - 0.17 \approx 48.90 \text{ N}$ \end{soln}

\bigskip

\example{Star $A$ of mass $6.0\times10^{30} \text{ kg}$ and star $B$ of mass $1.5\times10^{30} \text{ kg}$ are separated by a distance of $2.0 \times 10^{12} \text{ m}$. (a) What is the field strength at the mid-point $P$ of the two stars? (b) If a comet of mass $4.0\times10^6 \text{ kg}$ is at the mid-point, what force does it experience?}

\bigskip

\example{The earth has a radius of 6370 km. (a) Find the mass of the earth.
	\footnote{British scientist Henry Cavendish devised an experiment in 1798 to measure the gravitational force between masses in his laboratory. He was the first man to yield accurate values for the gravitational constant $G$. Then he was able to carry out this calculation, referred by himself as `weighing the world'.} 
(b) Find the acceleration of free fall at the top of Mount Everest. (height of Mount Everest $H\approx 8.8 \text{ km}$)}

\begin{soln}
    
 consider acceleration of free fall near surface of earth:
\begin{equation*}
	g_s = \frac{GM}{R^2} \RA 9.81 = \frac{6.67\times10^{-11}\times M}{(6.37\times10^6)^2} \RA M \approx 5.97\times10^{24} \text{ kg}
\end{equation*}
at top of Mount Everest:
\begin{equation*}
g_\text{ME} = \frac{GM}{(R+H)^2} = \frac{6.67\times10^{-11}\times 5.97\times10^{24}}{(6.37\times10^6+8.8\times10^3)^2} \approx 9.78 \text{ N kg}^{-1} 
\end{equation*}
\end{soln}

\bigskip

\example{field around the earth}

\bigskip

\example{field near earth's surface}

\bigskip

\example{A meteor is travelling towards a planet of mass $M$. When it is at a distance of $r_1$ from centre of $M$, it moves at speed $v_1$. When it is $r_2$ from $M$, it moves at speed $v_2$. Assume only gravitational force applies, establish a relationship between these quantities.}

\begin{soln}
    energy conservation: $\text{K.E.} + \text{G.P.E.} = \text{const}$ $\ra$ $\frac{1}{2}mv_1^2+\left(-\frac{GMm}{r_1}\right) = \frac{1}{2}mv_2^2+\left(-\frac{GMm}{r_2}\right)$ \end{soln}

\bigskip

\example{If an object is thrown from the surface of a planet at sufficiently high speed, it might escape from the influence of the planet's gravitational field. The minimum speed required is called the \emph{escape velocity}\index{escape velocity}. Using the data from previous examples, find the escape velocity from the surface of earth.}

\bigskip

\example{A star $A$ of mass $M_A=1.5\times10^{30}$ kg and a planet $B$ of mass $M_B=6.0\times10^{26}$ kg form an isolated astronomical system. Point $P$ is between $A$ and $B$, and is at distance $r_A=2.0\times10^{12}$ m from $A$, and distance $r_B=8.0\times10^{10}$ m from $B$. (a) Find the gravitational potential at $P$. (b) A meteor is initially at very large distance from the system with negligible speed. It then travels towards point $P$ due to the gravitational attraction. Find its speed when it reaches $P$.}

\begin{soln} gravitational potential at $P$: $\, \varphi_P = \varphi_A + \varphi_B = \left(-\frac{GM_A}{r_A}\right) + \left(-\frac{GM_B}{r_B}\right)$
\begin{equation*}
\varphi_P = -6.67\times10^{-11}\times\left(\frac{1.5\times10^{30}}{2.0\times10^{12}} + \frac{6.0\times10^{26}}{8.0\times10^{10}}\right) \approx -5.05\times10^{7} \text{ J kg}^{-1}
\end{equation*}

gain in K.E. = loss in G.P.E.: $\, \frac{1}{2}mv^2 = m\Delta \varphi \RA v^2 = 2(\varphi_\infty - \varphi_P)=-2\varphi_P$
\begin{equation*}
v=\sqrt{-2\times(-5.05\times10^7)} \approx 1.01\times10^4 \mps 
\end{equation*}\end{soln}

\bigskip

\example{The Moon may be considered to be an isolated sphere of radius $R=1.74 \times 10^3$ km. The gravitational potential at the surface of the moon is about $-2.82 \times 10^6 \text{ J kg}^{-1}$. (a) Find the mass of the moon. (b) A stone travels towards the moon such that its distance from the centre of the moon changes from $3R$ to $2R$. Determine the change in gravitational potential. (c) If the stone starts from rest, find its final speed.}

\begin{soln} at surface: $\varphi(R) = -\frac{GM}{R} \RA -2.82\times10^6 = -\frac{6.67\times10^{-11}\times M}{1.74\times10^6} \RA M = 7.36 \times 10^{22} \text{ kg}$

\eqyskip from $3R$ to $2R$: $\Delta \varphi = \varphi_{(3R)} - \varphi_{(2R)} = \left(-\frac{GM}{3R}\right) - \left(-\frac{GM}{2R}\right) = \frac{GM}{6R} = \frac{2.82\times10^6}{6} \approx 4.70 \times 10^5 \text{ J kg}^{-1}$

note this change is a \emph{decrease} in gravitational potential

gain in K.E. = loss in G.P.E.: $\, \frac{1}{2}mv^2 = m\Delta \varphi \RA v = \sqrt{2\Delta \varphi} = \sqrt{2\times4.70 \times 10^5} \approx 970 \mps$ \end{soln}

\bigskip

\section{Oscillations}

\example{Compare the two oscillations from the $x$-$t$ graph below.}

\bigskip

\example{A simple harmonic oscillator is displaced by 6.0 cm from its rest position and let go at $t=0$. Given that the period of this system is 0.80 s, state an equation for its displacement-time relation.}

\bigskip

\example{A simple harmonic oscillator is initially at rest. At $t=0$, it is given an initial speed in the negative direction. Given that the frequency is 1.5 Hz and the amplitude is 5.0 cm, state an equation for its displacement-time relation.}

\bigskip

\example{The motion of a simple pendulum is approximately simple harmonic. As the pendulum swings from one side to the other end, it moves through a distance of $6.0 \text{ cm}$ and the time taken is $1.0 \text{ s}$. (a) State the period and amplitude. (b) Find the greatest speed during the oscillation. (c) Find its speed when displacement $x=1.2\text{ cm}$.}

\begin{soln} period: $T=2\times 1.0 = 2.0 \text{ s}$, and amplitude: $x_0 = \frac{1}{2}\times6.0 = 3.0 \text{ cm}$

angular frequency: $\omega = \frac{2\pi}{T} = \frac{2\pi}{2.0} = \pi \radps$

greatest speed: $v_\tmax = \omega x_0 = \pi \times 3.0 \approx 9.4 \text{ cm s}^{-1}$

speed at $1.2 \text{ cm}$: $v = \omega\sqrt{x_0^2 - x^2} = \pi \times \sqrt{3.0^2-1.2^2} \approx 8.6 \text{ cm s}^{-1}$ \end{soln}

\bigskip

\example{Given the $x$-$t$ graph of a simple harmonic oscillator. (a) Find its speed at $t=0$. (b) Find its greatest speed. (b) Find its acceleration at $t=1.0$ s.}

\bigskip

\example{A block of mass 150 g at the end of a spring oscillates with a period of 0.80 s. The maximum displacement from its rest position is 12 cm. Find the energy of the vibration.}

\begin{soln}\begin{equation*}
	E = \frac{1}{2}m\omega^2x_0^2 = \frac{1}{2}m\left(\frac{2\pi}{T}\right)^2x_0^2 = \frac{1}{2} \times 0.15 \times \frac{4\pi^2}{0.80^2}\times0.12^2 \approx 6.7\times10^{-2} \text{ J} 
\end{equation*}
\end{soln}

\bigskip

\example{An oscillator is composed of a block of mass $m=250$ g and a spring of $k=1.6$ N/cm. It is displaced by 5.0 cm from its rest position and set free. (a) What is its angular frequency? (b) what is the initial vibrational energy? (c) After a few oscillations, 40\% of its energy is lost due to damping. What is its new amplitude?}

\begin{soln}
     angular frequency: $\omega = \sqrt{\frac{k}{m}} = \sqrt{\frac{160}{0.25}} \approx 25.3 \radps \,$

energy of oscillator: $E=\frac{1}{2}m\omega^2x_0^2 = \frac{1}{2} \times 0.25 \times 25.3^2 \times 0.050^2 = 0.20 \text{ J} $

An easier approach: $E = \frac{1}{2}kx_0^2 = \frac{1}{2} \times 160 \times 0.050^2 = 0.20 \text{ J}$.

since $E \propto x_0^2$, so: $\frac{E'}{E} = \frac{x^{\prime2}_0}{x_0^2} \RA  60\% = \frac{x^{\prime2}_0}{x_0^2} \RA x'_0 = \sqrt{0.6} x_0 = \sqrt{0.6} \times 5.0 \approx 3.9$ cm \end{soln}

\bigskip

\section{Ideal Gases}

\example{Find the number of molecules in 160 grams of argon-40 gas.}

\begin{soln} Amount of gas: $n=\frac{m}{M} = \frac{160 \text{ g}}{40 \text{ g mol}^{-1}} = 4.0 \text{ mol}$

number of gas molecules: $N = n N_A = 4.0 \text{ mol} \times 6.02\times10^{23} \text{ mol}^{-1} \approx 2.41 \times 10^{24}$ \end{soln}

\bigskip

\example{If a gas is heated with its volume fixed, how does the pressure change?}

\begin{soln}At higher temperature, gas molecules move faster

they will collide \emph{harder} and produce a greater force upon each collision

they will also collide more \emph{frequently} with the container

so pressure of the gas will increase \end{soln}

\bigskip

\example{A sealed cylinder of volume of 0.050 m$^3$ contains 75 g of air. The molar mass of air is 29 g mol$^{-1}$. (a) Find the air pressure when its temperature is $30^\circ$C. (b) The gas is allowed to expand with its pressure fixed. Find the temperature of the gas when the volume doubles.}

\begin{soln}
    
amount of gas: $n=\frac{m}{M} = \frac{75}{29} \approx 2.59 \text{ mol}$

pressure at $30^\circ$C: $p = \frac{nRT_1}{V_1} = \frac{2.59\times8.31\times(30+273)}{0.050} \approx 1.30\times10^5 \text{ Pa}$


pressure fixed, so $V \propto T \RA \frac{T_2}{T_1} = \frac{V_2}{V_1} = 2 \RA T_2 = 2\times (30+273) = 606 \text{ K} = 333^\circ\text{C}$ \end{soln}

\bigskip

\example{A gas cylinder holding 5000 cm$^3$ of air at a temperature of 27 $^\circ$C and a pressure of $6.0\times10^5 \text{ Pa}$ is used to fill balloons. Each balloon contains 1000 cm$^3$ of air at 27 $^\circ$C and $1.0\times10^5$ Pa when filled. (a) Find the initial amount of gas in the cylinder. (b) Find the number of balloons that can be filled.}

\begin{soln} initial amount of gas in cylinder: $n_0 = \frac{p_0 V}{RT} = \frac{9.0\times10^5\times5000\times10^{-6}}{8.31\times(27+273)} \approx 1.203 \text{ mol}$

final amount of gas in cylinder: $n_\text{remain} = \frac{pV}{RT} = \frac{1.0\times10^5\times5000\times10^{-6}}{8.31\times(27+273)} \approx 0.201 \text{ mol}$

Air will leave the cylinder to fill balloons only if pressure inside the cylinder is higher than pressure of the balloon. When the two pressures become equal, no more balloons can be filled, there will be some air remain in cylinder...

amount of gas in each balloon: $n_\text{b} = \frac{pV_b}{RT} = \frac{1.0\times10^5\times1000\times10^{-6}}{8.31\times(27+273)} \approx 0.040 \text{ mol}$

number of balloons: $N = \frac{n_0 - n_\text{remain}}{n_\text{b}} = \frac{1.203-0.201}{0.040} \approx 25$ \end{soln}

\bigskip

\example{A storage cylinder has a volume of $5.0\times10^{-4}\text{ m}^3$. The gas is at a temperature
	of 300 K and a pressure of $4.0 \times 10^6$ Pa.	(a) Find the number of molecules in the cylinder. (b) The gas molecules slowly leak from the cylinder at a rate of $1.6 \times 10^{16} \text{ s}^{-1}$. Find the time, in days, after which the pressure will reduce by 5.0\%.}

\begin{soln} initial number of molecules: $N_0 = \frac{p_0V}{kT} = \frac{4.0\times10^6\times 5.0\times10^{-4}}{1.38\times10^{-23}\times300}\approx 4.83\times10^{23}$

\eqyskip

volume fixed, so $N \propto p \RA \frac{\Delta N}{N_0} = \frac{\Delta p}{p_0} = 5.0\%$

number of molecules escaped: $\Delta N = 0.05\times4.83\times10^{23} \approx 2.42 \times10^{22}$

time needed: $t = \frac{2.42 \times10^{22}}{1.6 \times 10^{16}} \approx 1.51 \times 10^6 \text{ s} \approx 17.4 \text{ days}$ \end{soln}

\bigskip

\example{A mass of 20 g helium-4 at a temperature of 37$^\circ$C has a pressure of $1.2\times10^5 \text{ Pa}$. Each helium-4 atom has a diameter of 280 pm. (a) Find the volume occupied by the gas and the volume of atoms in this gas. (b) Compare the two volumes, suggest whether this gas can be considered as an ideal gas.}

\begin{soln}
    
number of helium molecules: $N = nN_A = \frac{m}{M} \times N_A = \frac{20}{4.0} \times 6.02\times10^{23} \approx 3.01\times10^{24} $

\eqyskip

volume of gas: $V_\text{gas} = \frac{NkT}{p} = \frac{3.01\times10^{24}\times1.38\times10^{-23}\times (37+273)}{1.2\times10^5} \approx 0.107 \text{ m}^3$

\eqyskip

volume of one atom: $V_\text{atom} = \frac{4}{3}\pi r^3 = \frac{4}{3} \pi\times(140\times10^{-12})^3 \approx 1.15\times10^{-29} \text{ m}^3$

volume of all atoms: $V_\text{atoms} = N V_\text{atom} = 3.01\times10^{24} \times 1.15\times10^{-29} \text{ m}^3 \approx 3.46 \times 10^{-5} \text{ m}^{3}$

$V_\text{gas} \gg V_\text{atoms}$, so this gas can approximate to an ideal gas \end{soln}

\bigskip

\example{Air consists of oxygen (O$_2$, molar mass $32\text{ g mol}^{-1}$) and nitrogen (N$_2$, molar mass $28\text{ g mol}^{-1}$). (a) Calculate the mean translational kinetic energy of these molecules at 300 K. (b) Estimate the typical speed for each type of the molecule.}

\begin{soln} mean K.E. of single molecule: $\avg{E_k} = \frac{3}{2}kT = \frac{3}{2} \times 1.38\times10^{-23} \times 300 \approx 6.21\times 10^{-21} \text{ J}$

{
	
	\centering
	
	$\avg{E_k} = \frac{1}{2}m\avg{v^2} = \frac{3}{2} kT \RA \frac{1}{2} \frac{M}{N_A} \avg{v^2} = \frac{3}{2} kT \RA \avg{v^2} = \frac{3kN_AT}{M} = \frac{3RT}{M}$
	
}

\eqyskip

for oxygen molecule: $v_\text{O$_2$} \approx  \sqrt{\frac{3\times8.31\times300}{0.032}} \approx 483 \mps$

\eqyskip

for nitrogen molecule: $v_\text{N$_2$} \approx  \sqrt{\frac{3\times8.31\times300}{0.028}} \approx 517 \mps$ \end{soln}

\bigskip

\example{A cylinder container initially holds a gas of helium-4 at a temperature of 54\OC. (a) Find the mean square speed of these helium atoms. (b) If the temperature is raised to 540\OC, find the r.m.s. speed of the atoms.}

\begin{soln} mass of one helium-4 atom: $m=4\text{u} = 4\times 1.66\times10^{-27} \approx 6.64\times10^{-27} \text{ kg}$

at 54\OC: $\, \frac{1}{2}m\avg{v^2} = \frac{3}{2}kT \RA \avg{v^2} = \frac{3kT}{m} = \frac{3\times1.38\times10^{-23}\times(54+273)}{6.64\times10^{-27}} \approx 2.04 \times 10^6 \text{ m}^2 \text{ s}^{-2}$

\eqyskip 

note relation between $v$ and $T$: $\, \avg{v^2} \propto T \RA \frac{\avg{v^{\prime 2}}}{\avg{v^2}} = \frac{T'}{T} \RA v'_\text{rms} = \sqrt{\frac{T'}{T}}\times v_\text{rms}$

\eqyskip 

at 540\OC: $\, v'_\text{rms} = \sqrt{\frac{540+273}{54+273}} \times \sqrt{2.04 \times 10^6} \approx 2.25 \times 10^3 \mps$ \end{soln}

\bigskip

\section{Thermodynamics}

\example{A 3.0 kW electric kettle contains 0.5 kg of water already at its boiling point. Neglecting heat losses, determine how long it takes to boil dry. ($L_\text{water} = 2.26 \times 10^6 \text{ J kg}^{-1}$)}

\bigskip

\example{A student measures specific latent heat of fusion for ice. He uses an electric heater to melt ice but the insulation is not perfect. The experiment is carried out twice, with the heater operating at different powers. Use the data table to calculate specific latent heat of fusion.}

\bigskip

\example{A block of 30 g ice at $-20^\circ$C is added to a large cup of 270 g water at 80$^\circ$C. Assume there is no energy lost, what is the final temperature of the mixture? (data: specific heat capacity of water is 4200 J kg$^{-1}$ K$^{-1}$, specific heat capacity of ice is 2100 J kg$^{-1}$ K$^{-1}$, specific latent heat of ice is $3.3\times10^5$ J kg$^{-1}$.}

\bigskip

\example{A 1.00 kg aluminium block is heated using an electrical heater. The current in the heater is 4.2 A and the p.d. across is 12 V. Measurements of the rising temperature are represented by the graph. Determine specific heat capacity  of aluminium.}

\bigskip

\example{A gas is heated by supplying it with 25 kJ of energy. The gas expands so that the volume increases by 0.10 m$^3$. Assume the gas has a fixed pressure of 150 kPa during the process. Calculate the change in internal energy.}

\begin{soln} amount of work done: $W = p \Delta V = 150 \times 0.10 = 15 \text{ kJ}$

but gas expands means work is done against surroundings, so this is negative work

change in internal energy: $\Delta U = Q +  W = (+25) + (-15) = +10 \text{ kJ}$ \end{soln}

\bigskip

\example{Use the idea of internal energy and the first law of thermodynamics, explain why boiling water requires heat supply.}

\begin{soln} boiling occurs at constant temperature, so $\Delta E_k = 0$

but separation between molecules increased, so $\Delta E_p >0$

by definition, internal energy $U=E_k+E_p$, so $\Delta U >0$

during boiling, there is an increase in volume, so work against surroundings, $W<0$

recall first law of thermodynamics $\Delta U = Q + W$, must have $Q>0$

this means heat must be supplied for boiling processes \end{soln}

\bigskip

\example{When you pump up a bicycle tyre, the temperature of air inside the tyre will go up. Explain why this happens using the first law of thermodynamics.}

\begin{soln} pumping up tyre involves compressing gas, so positive work is done: $W>0$

for each stroke, there is little time for heat transfer, so $Q \approx 0$

according to first law of thermodynamics $\Delta U = Q + W \RA \Delta U >0$

by definition, internal energy $U=E_k + E_p \RA \Delta U = \Delta E_k + \Delta E_p>0$

but for gas, there is negligible intermolecular force, so $\Delta E_p=0$, then must have $\Delta E_k>0$

K.E. of molecules is proportional to temperature, higher K.E. so higher temperature \end{soln}

\bigskip

\example{An ideal gas of 0.080 mol is initially at state $A$ and then undergoes a cycle $ABCA$. The variation of its pressure $p$ with its volume $V$ is shown on the graph.}

\bigskip

\section{Electrostatics}

\example{The hydrogen atom has a radius of about 53 pm. Estimate the electric force between the proton and the orbiting electron.}

\begin{soln}\begin{equation*}
	F = \frac{Q_p Q_e}{\ec r^2} = 8.99\times10^9 \times \frac{(1.60\times10^{-19})^2}{(53\times10^{-12})^2} \approx 8.2 \times 10^{-8} \text{ N} 
\end{equation*}
\end{soln}

\bigskip

\example{A \emph{Van de Graaff generator} produces sparks when its surface electric field strength $4.0 \times 10^4 \Vpcm$. If the diameter of the sphere is $40 \text{ cm}$, what is the charge on it?}

\bigskip

\example{Two identical metal spheres of radius $20 \text{ cm}$ carry charges $+2.0 \text{ $\mu$C}$ and $-1.0 \text{ $\mu$C}$ respectively. There is a $60 \text{ cm}$ gap between them. (a) Find the electric field strength midway along the line joining their centres. (b) A dust particle carrying a charge of $-1.3\times10^{-8} \text{ C}$ is at this position. Find the electric force it experiences.}

\bigskip

\example{Sketch the electric field around a positively-charged sphere or a negatively charged sphere, and explain why they can be considered as point charges.}

\bigskip

\example{Field lines between two oppositely-charged large metal plates.}

\bigskip

\example{Field pattern due to two charges of equal magnitude.}

\bigskip

\example{To gain information about the gold nucleus ($_{\phantom{0}79}^{197}\text{Au}$), we fire $\alpha$-particles ($_2^4\alpha$) towards a thin gold foil. The size of a typical nucleus is about $10^{-14}$m, what is the minimum initial speed for $\alpha$-particles so that radius of gold nucleus can be determined?}

\bigskip

\example{An electron is accelerated from rest through a potential difference of 600 V.
	
Find the final speed of the electron.}

\begin{soln} gain in K.E. = change in electric P.E. $\RA \frac{1}{2}mv^2 = q\Delta V$
\eqyskip\begin{equation*}
	v = \sqrt{\frac{2q\Delta V}{m}} = \sqrt{\frac{2\time 1.60\times10^{-19}\times600}{9.11\times10^{-31}}} \approx 1.45\times10^7 \mps 
\end{equation*}
\end{soln}

\bigskip

\example{Two small metal spheres $A$ and $B$ are in a vacuum. Sphere $A$ has charge $+20$ pC and sphere $B$ has charge $+84$ pC. The arrangement is shown below.}

\bigskip

\example{$A$ and $B$ are two positively-charged spheres of radius 1.0 cm. A proton $P$ initially at rest on the surface of $A$ moves along the line joining the centres of the two spheres. The variation with distance $x$ from the centre of $A$ of electric potential $V$ at point $P$ is given.}

\bigskip

\example{Field lines and equipotential lines due to two charges of various magnitudes - you might be given an example and asked to sketch the lines and equipotentials.}

\bigskip

\example{Consider the electric field due to a metal sphere of radius $R$ carrying charge $Q$. Plot the variation with the distance $r$ from sphere's centre of the field strength, and the variation with $r$ of the electric potential.}

\bigskip

\example{The variation of electric potential near a charged object is shown on the graph. Calculate the electric field strength at 5.0 cm from the centre of the object.}

\bigskip

\example{The variation of electrical potential along a certain line is shown. State and explain where in the field an electron will experience the greatest force.}

\bigskip

\example{electric field due to an isolated point charge}

\bigskip

\section{Capacitors}

\example{When the p.d. across a capacitor of $1.8\times10^{-4} \text{ F}$ is increased from 10 V to 20 V, how much additional energy is stored?}

\begin{soln}energy change: $\Delta W = W_f - W_i = \frac{1}{2}CV_f^2 - \frac{1}{2}CV_i^2 = \frac{1}{2}\times 1.8\times10^{-4}\times(20^2-10^2) = 2.7\times10^{-2} \text{ J}$ \end{soln}

\bigskip

\example{A capacitor with capacitance $C_0$ is charged to a p.d. $V_0$. It is disconnected from the power supply, and then connected across an identical capacitor. Discuss the change in p.d., and change in energy stored in the system.}

\begin{soln} Initial charge $Q=C_0 V_0$, initial energy stored $W_0 = \frac{1}{2} C V_0^2$

combined capacitance: $C= C_0 + C_0 = 2C_0$

charge is conserved, so final p.d. across: $V = \frac{Q}{C} = \frac{C_0 V_0}{2C_0} = \frac{1}{2} V_0$


charge shared between capacitors, the first capacitor loses half its charge and p.d.

final energy stored in system: $W = \frac{1}{2}CV^2 = \frac{1}{2} \times 2C_0 \times (\frac{1}{2}V_0)^2 = \frac{1}{4} C_0 V_0^2 \ra W = \frac{1}{2} W_0$

half of stored energy is lost as heat when electrons flow between two capacitors \end{soln}

\bigskip

\example{For the circuit shown below, find the p.d. across each of the capacitor.}

\bigskip

\example{$C_1=C_2=C_3=C_4=10\muF$, calculate the capacitance of the network (a) and (b).}

\bigskip

\example{Four identical capacitors are arranged as shown. Each capacitor can withstand a maximum p.d. of 12 V, what is the maximum safe p.d. to be applied between the terminals?}

\bigskip

\section{Magnetic Fields}

\example{The magnetic field around a bar magnet}

\bigskip

\example{Determine the direction of the magnetic force acting. Check for yourself.}

\bigskip

\example{A wire of 0.80 m carrying a current of 5.0 A is lying at right angles to a magnetic field, it experiences a force of 0.60 N, what is the flux density?}

\bigskip

\example{Torque on a rectangular metal frame in uniform magnetic field}

\bigskip

\example{A U-shaped magnet is placed on a balance with a wire suspended above it as shown. Magnetic flux density between the poles is about 0.50 T. The part of the wire that is in the field is of length 8.0 cm. The balance initially shows a reading of 500.0 g. When a current of 10 A flowing into the paper is switched on, what is the new reading on the balance?}

\bigskip

\example{Use the left-hand rule to find the direction of the magnetic forces acting on the following moving charges. Check yourself.}

\bigskip

\example{An $\alpha$-particle travelling at $2.5\times10^4 \mps$ enters a region of uniform magnetic field. The field has flux density of 5.4 mT and is normal to direction of particle's velocity. What is the radius of $\alpha$-particle's path?}

\begin{soln} Radius of circular arc: $r=\frac{mv}{Bq} = \frac{4\times 1.66\times10^{-27}\times2.5\times10^4}{5.4\times10^{-3}\times2\times1.60\times10^{-19}} \approx 0.096 \text{ m} \approx 9.6 \text{ cm}$ \end{soln}

\bigskip

\section{Electromagnetic Induction}

\example{A coil of 80 turns is wound tightly around a solenoid with a cross-sectional area of 35 cm$^2$. The flux density at the centre of the solenoid is 75 mT. (a) What is the flux linkage in the coil? (b) The current in the solenoid is \emph{reversed} in direction in a time of 0.40 s, what is the average e.m.f. induced?}

\begin{soln} flux linkage: $\Phi = NBA = 80 \times 75\times10^{-2} \times 35 \times 10^{-4} = 0.21 \text{ Wb}$

average e.m.f. induced: $\mathcal{E} = \frac{\Delta \Phi}{\Delta t} = \frac{(+\Phi)-(-\Phi)}{\Delta t} = \frac{2\times 0.21}{0.40} = 1.05 \text{ V}$ \end{soln}

\bigskip

\example{A coil of 200 turns is connected to a resistor $R=4.0$ $\Omega$. The coil has a diameter 10 cm. Initially a bar magnet is at great distance from the coil. The magnet is then inserted into the coil and field inside coil becomes 0.40 T. The process occurs within a duration of 2.0 s.}

\bigskip

\example{Two parallel metal tracks are separated by a distance of $L=45$ cm and placed in a uniform magnetic field of 0.80 T. The tracks are connected to a resistor $R=6.0$ $\Omega$. A long metal rod is being pushed under an external force at $v=3.0 \mps$ along the tracks as shown.}

\bigskip

\section{Alternating Currents}

\example{The variation with time of an alternating current in a resistor of $120 \Omega$ is shown. (a) What is the value of its r.m.s. current? (b) What is the mean power dissipated by the resistor?}

\bigskip

\example{An a.c. power supply produces a sinusoidal output across a resistor of 30 $\Omega$. The maximum voltage is found to be 75 V. Find energy dissipated in the resistor in 2.0 minutes.}

\bigskip

\example{An oscilloscope displays an a.c. voltage signal as shown.}

\bigskip

\example{An ideal transformer has 200 turns on the primary coil and 5000 turns on the secondary	coil. The r.m.s. input voltage to the primary coil is 8.0 V. What is the \emph{peak} voltage across a resistor connected to the secondary coil?}

\begin{soln}\begin{equation*}
\frac{V_{s,0}}{V_{p,0}} = \frac{N_s}{N_p} \RA 
\frac{V_{s,0}}{\sqrt{2}V_{p,\text{rms}}} = \frac{N_s}{N_p} \RA
\frac{V_{s,0}}{8.0\times\sqrt{2}} = \RA
V_{s,0} \approx 2830 \text{ V}   
\end{equation*}
\end{soln}

\bigskip

\section{Quantum Physics Basics}

\example{A laser emits red light of 650 nm at a power rating of 2.0 mW. (a) What is the energy carried by one photon? (b) How many photons are emitted per second?}

\bigskip

\example{Given that work function energy of gold is 4.9 eV. Find the longest wavelength of electromagnetic wave that could release electrons from gold.}

\begin{soln}
    
 threshold frequency: $f_0 = \frac{\Phi}{h} = \frac{4.9\times1.60\times10^{-19}}{6.63\times10^{-34}} \approx 1.18\times10^{18} \text{ Hz}$

threshold wavelength: $\lambda_0 = \frac{c}{f_0} = \frac{3.00\times10^8}{1.18\times10^{15}} \approx 2.54\times10^{-7} \text{ m} \,$ (ultraviolet light) \end{soln}

\bigskip

\example{A beam of light has wavelength 600 nm, cross-sectional area $0.16 \text{ cm}^2$ and power 5.0 mW. The beam is normally incident onto a surface and is completely absorbed. Calculate, for a time of 1.0 s, (a) the number of photons incident onto the surface, (b) the change of total momentum of the photons, (c) the light pressure on the surface.}

\begin{soln}
    
 energy of one photon: $E= \frac{hc}{\lambda} = \frac{6.63\times10^{-34} \times 3.0\times10^8}{600\times10^{-9}} \approx 3.32 \times 10^{-19} \text{ J}$ 1.5083


number of photons arriving in 1.0 s: $N = \frac{5.0\times10^{-3} \times 1.0}{3.32 \times 10^{-19}} \approx 1.51\times10^{16}$

momentum for one photon: $p = \frac{h}{\lambda} = \frac{6.63\times10^{-34}}{600\times10^{-9}} \approx 1.11 \times 10^{-27} \text{ kg m s}^{-1}$

change of total momentum: $\Delta P = Np = 1.51\times10^{16} \times 1.11 \times 10^{-27} \approx 1.67 \times 10^{-11} \text{ kg m s}^{-1}$

average force due to these photons: $F = \frac{\Delta P}{\Delta t} = \frac{1.67 \times 10^{-11}}{1.0} \approx 1.67 \times 10^{-11} \text{ N}$

light pressure on surface: $p = \frac{F}{A} = \frac{1.67 \times 10^{-11}}{0.16 \times 10^{-4}} \approx 1.04 \times 10^{-6} \text{ Pa}$ \end{soln}

\bigskip

\example{What is the wavelength of a human of 70 kg walking at around $2.0\mps$?}

\begin{soln}
    $\lambda = \frac{h}{mv} = \frac{6.63\times10^{-34}}{75\times 2} \approx 4.7\times 10^{-36} \text{ m}$

this wavelength is too small compared with any obstacle we encounter in everyday lives

so human bodies do not exhibit noticeable wave behaviour \end{soln}

\bigskip

\example{An electron is accelerated from rest through a potential difference of 50 V. What is the de Broglie wavelength of this electron?}

\begin{soln}
K.E. of electron equals change in electric P.E., so
$\frac{1}{2}mv^2 = qV \RA v = \sqrt{\frac{2qV}{m}} = \sqrt{\frac{2\times1.60\times 10^{-19} \times 50}{9.11\times10^{-31}}} \approx 4.19\times10^6 \mps$


wavelength of electron: $\lambda = \frac{h}{mv} = \frac{6.63\times10^{-34}}{9.11\times10^{-31}\times 4.19\times10^6} \approx 1.74 \times10^{-10} \text{ m}$

this wavelength is comparable to scale of atomic spacing (also around $10^{-10}\text{ m}$)

so these electron can be \emph{diffracted} by solid crystals\end{soln}

\bigskip

\example{Some electron energy levels of the hydrogen atom is shown. Find the longest and the shortest wavelength produced by electron transitions between the energy levels given.}

\bigskip

\example{The emission spectrum of the hydrogen atom consists of a number of wavelengths in the visible spectrum. Given that the visible spectrum consists of light of wavelengths within the range from 380 nm to 740 nm, also use data in Example \ref{ex-hydrogen}, find the energies of visible photons that can be produced by transitions between the energy levels shown.}

\bigskip

\example{A white light is incident on a cloud of cool hydrogen gas. In the emergent spectrum, a dark line is observed at a wavelength of 435 nm. Determine the energy change that gives rise to this dark line.}

\begin{soln}
    
 photon absorbed: $E = \frac{hc}{\lambda} = \frac{6.63\times10^{-34}\times 3.0 \times10^8}{435\times10^{-9}} \approx 4.57\times10^{-19} \text{ J} = \frac{4.57\times10^{-19}}{1.60\times10^{-19}} \text{ eV} \approx 2.86 \text{ eV}$

note that $3.40 - 0.54 = 2.86 \text{ eV}$, so this dark line is due to the electron transition: $E_2 \to E_5$ \end{soln}

\bigskip

\example{Silicon has a band gap $E_g \approx 1.1 \text{ eV}$. State and explain whether it becomes conducting when exposed to red light of $600 \text{ nm}$.}

\begin{soln}
    
 Energy of radiation: $E=\frac{hc}{\lambda} = \frac{6.63\times10^{-34}\times3.0\times10^8}{600\times10^{-9}} \approx 3.32\times10^{-19} \text{ J} \approx 2.1 \text{ eV} > E_g$

this energy is sufficient to make valence electrons in silicon to cross band gap

silicon now possesses conducting electrons and holes therefore becomes conducting \end{soln}

\bigskip

\section{Nuclear Physics}

\example{When an antiproton (antiparticle of a proton) collides with a proton, they are annihilated and two photons of equal energy are formed. (a) What is the energy of each photon? (b) What is the photon frequency?}

\begin{soln}Rest energy of proton becomes photon energy:

{

\centering

$E_\gamma=m_pc^2 = 1.67\times10^{-27} \times (3.0\times10^8)^2 \approx 1.50\times10^{-10} \text{ J}$

}

photon frequency: $f=\frac{E_\gamma}{h} = \frac{1.50\times10^{-10}}{6.63\times10^{-34}} \approx 2.27 \times10^{23} \text{ Hz}\,$ (high-frequency $\gamma$-photon) \end{soln}

\bigskip

\example{An iron-56 nucleus ($^{56}_{26}$Fe) has a mass of $9.288 \times 10^{-26} \text{ kg}$. Calculate the nuclear binding energy per nucleon, in MeV, for $^{56}_{26}$Fe.}

\begin{soln} mass defect: $\Delta m = 26m_\text{p} + (56-26)m_\text{n} - m_\text{Fe} = 26\times 1.673\times10^{-27} + 30 \times1.675\times10^{-27} - 9.288 \times 10^{-26}$

so we find $\Delta m = 8.68 \times 10^{-28} \text{ kg}$

binding energy: $E_B = \Delta m c^2 = 8.68 \times 10^{-28} \times (3.00\times10^8)^2 \approx 7.812 \times 10^{-11} \text{ J}$

binding energy per nucleon: $\epsilon_b = \frac{E_b}{A} = \frac{7.812 \times 10^{-11}}{56} \approx 1.395 \times 10^{-12} \text{ J}$

\eqyskip convert into MeV: $\epsilon_b = \frac{1.395 \times 10^{-12}}{1.60\times10^{-13}} \text{ MeV} \approx 8.72 \text{ MeV}$ \end{soln}

\bigskip

\example{The sun is a huge nuclear reactor. It fuses hydrogen nuclei into helium nuclei to produce large amounts of energy. The reaction can be expressed by the formula:}

\bigskip

\example{Uranium-235 nuclei ($^{235}_{\phantom{2}92}$U) bombarded with slow neutrons can undergo nuclear reaction:}

\bigskip

\example{A sample containing $6.0\times10^7$ iodine-131 nuclei has an activity of 60 Bq. (a) What is the decay constant of iodine-131? (b) How many nuclei remain undecayed after 20 days?}

\begin{soln} decay constant: $\lambda = \frac{A}{N} = \frac{60}{6.0\times10^7} \approx 1.0 \times 10^{-6} \text{ s}^{-1}$

number of remaining nuclei: $N=N_0 \mathrm{e}^{-\lambda t} = 6.0\times10^7 \times \mathrm{e}^{-1.0 \times 10^{-6} \times20\time24\times3600} \approx 1.1\times10^7$ \end{soln}

\bigskip

\example{$^{24}_{11}\text{Na}$, an isotope of sodium, has a decay constant of $1.28\times10^{-5} \text{ s}^{-1}$. Suppose a sample initially contains a mass of 9.0 $\mu\text{g}$ of $^{24}_{11}\text{Na}$. Find (a) initial number of $^{24}_{11}\text{Na}$ nuclei, (b) initial activity of this sample, (c) the activity after 48 hours.}

\begin{soln} initial number of nuclei: $N_0 = \frac{9.0 \text{ }\mu\text{g}}{24\text{u}} =  \frac{9.0\times10^{-9}}{24\times1.66\times10^{-27}} \approx 2.26\times10^{17}$

initial acitivy: $A_0 = \lambda N_0 = 1.28\times10^{-5} \times 2.26\times10^{17} \approx 2.89\times10^{12} \text{ Bq}$

activity after 20 hours: $A=A_0\mathrm{e}^{-\lambda t} = 2.89\times10^{12} \times \mathrm{e}^{-1.28\times10^{-5} \times48\times3600} \approx 3.17\times10^{11} \text{ Bq}$ \end{soln}

\bigskip

\example{Radium-224 has a half-life of 3.63 days. (a) What is the activity from 6.0 mg of pure radium-224? (b) How many radium-224 nuclei have undergone decay after 10 days?}

\begin{soln}
    
 The decay constant: $\lambda = \frac{\ln 2}{\halflife} = \frac{\ln 2}{3.63\times24\times3600} \approx 2.21\times10^{-6} \text{ s}^{-1}$

\eqyskip initial number of Ra-224 nuclei: $N_0 = \frac{6.0 \text{ mg}}{224 \text{ u}} = \frac{6.0\times10^{-6}}{224\times1.66\times10^{-27}} \approx 1.61\times10^{19}$

initial acitivy: $A_0 = \lambda N_0 = 2.21\times10^{-6} \times 1.61\times10^{19} \approx 3.57\times10^{13} \text{ Bq}$

number of undecayed nuclei: $N=N_0\mathrm{e}^{-\lambda t} = 1.61\times10^{19} \times \mathrm{e}^{-2.21\times10^{-6}\times10\times24\times3600} \approx 2.39\times10^{18}$

number of nuclei that have decayed: $\Delta N = N_0 - N = 1.61\times10^{19} - 2.39\times10^{18} \approx 1.37\times10^{19}$ \end{soln}

\bigskip

\example{Living trees contain a certain percentage of $^{14}_{\phantom{1}6}\text{C}$, an radioactive isotope of carbon that has a half-life of 5570 years. A sample of dead wood is found to have an activity of 0.42 Bq, while an equal mass of living wood has an activity of 1.60 Bq. Find the age of the dead wood.}

\begin{soln}\begin{equation*}
	A=A_0\mathrm{e}^{-\lambda t} = A_0\mathrm{e}^{ -\frac{\ln 2}{\halflife} t} \ra 0.42 = 1.60 \mathrm{e}^{ -\frac{\ln 2}{5570} t} 
 \end{equation*}
 \begin{equation*}
 \ra -\frac{\ln 2}{5570} t = \ln\left(\frac{0.42}{1.60}\right) \ra t\approx 10700 \text{ years} 
\end{equation*}\end{soln}

\bigskip

\example{A GM tube is placed close to a source of radioactive isotope. The variation with time $t$ of the measured count rate $R$ is shown. Determine the half-life of this isotope.}

\bigskip

\section{Electronics}

\example{For the sensing circuit shown, suggest how the output voltage changes as light level increases.}

\bigskip

\end{document}